\documentclass[11pt]{article}
\usepackage[T1]{fontenc}
\usepackage{lmodern}
\usepackage{microtype}
\usepackage[margin=1in]{geometry}
\usepackage{amsmath,amssymb,amsthm,mathtools}
\usepackage{mathrsfs}
\usepackage{booktabs,array,longtable}
\usepackage{enumitem}
\usepackage{tikz-cd}
\usepackage[hidelinks]{hyperref}
\usepackage[nameinlink,noabbrev]{cleveref}
\usepackage{xcolor}
\usepackage{needspace}
\hypersetup{pdftitle={Komorebi Seeded Two-Axis Calculus v0.8.0},
 pdfauthor={Komorebi P.},pdfsubject={Arithmetic resolution towers, completion defects, and effective lattice reconstruction}}
\setlength{\emergencystretch}{2em}
\setcounter{tocdepth}{2}

\setlist[itemize]{leftmargin=2em,itemsep=2pt,topsep=4pt}
\setlist[enumerate]{leftmargin=2.2em,itemsep=2pt,topsep=4pt}
\allowdisplaybreaks

\newtheorem{theorem}{Theorem}[section]
\newtheorem{proposition}[theorem]{Proposition}
\newtheorem{lemma}[theorem]{Lemma}
\newtheorem{corollary}[theorem]{Corollary}
\theoremstyle{definition}
\newtheorem{definition}[theorem]{Definition}
\newtheorem{axiom}[theorem]{Axiom}
\newtheorem{construction}[theorem]{Construction}
\newtheorem{example}[theorem]{Example}
\newtheorem{conjecture}[theorem]{Conjecture}
\newtheorem{problem}[theorem]{Open Problem}
\theoremstyle{remark}
\newtheorem{remark}[theorem]{Remark}

\newcommand{\K}{\mathbf K}
\newcommand{\Rcat}{\mathbf R}
\newcommand{\Tcat}{\mathbf T}
\newcommand{\Cat}{\mathbf{Cat}}
\newcommand{\Ob}{\operatorname{Ob}}
\newcommand{\Hom}{\operatorname{Hom}}
\newcommand{\Lift}{\operatorname{Lift}}
\newcommand{\Cofib}{\operatorname{Cofib}}
\newcommand{\Recon}{\operatorname{Recon}}
\newcommand{\Obs}{\operatorname{Obs}}
\newcommand{\Zoo}{\mathcal Z}
\newcommand{\Seeds}{\Sigma}
\newcommand{\Res}{\mathsf{Res}}
\newcommand{\Trans}{\mathsf{Trans}}
\newcommand{\id}{\operatorname{id}}
\newcommand{\colim}{\operatorname*{colim}}
\newcommand{\status}[1]{\textbf{Status: #1.}}

\title{\textbf{Komorebi Seeded Two-Axis Calculus}\\
\large Transformation, Resolution, Residual Generation, and Reconstruction\\[0.4em]
\normalsize Working Formalization v0.8.0}
\author{Komorebi P.}
\date{27 September 2026}

\begin{document}
\maketitle

\begin{abstract}
This manuscript develops a seed-generated calculus of transformations within a resolution, coarsening between resolutions, and the residuals of their exchange maps. Its categorical core uses indexed categories, coherent event diagrams, stable cofibers, and constructor-compatible Verdier quotients. It distinguishes a normalization endpoint, its comparison map, and the history-sensitive filtration of its residual. Excisive event diagrams reconstruct from principal ideals, while integral scalar and matrix models identify precise repair kernels and collapse conditions.

Version 0.8.0 adds genuinely nonidentity arithmetic coarsenings by derived reduction modulo powers of a squarefree integer. Coherent weighted exchange cells have explicit Koszul residuals. At a finite prime-power level, forcing an interaction of two nonunit scalar cells to be excisive collapses the whole perfect category. With several primes, the exact quotient removes precisely the selected prime components and commutes with finite coarsening. Uniform based-complex towers reconstruct coefficientwise, but repair need not commute with completion: the comparison's fiber consists of explicitly identified generic fibers.

An arithmetic fracture square retains the localized and completed descriptions together with their common generic data. A further theorem constructs a unique integral lattice from independently supplied localized modules, local lattices, and compatible generic identifications. Compatible operators descend to that lattice. These are restricted arithmetic reconstruction results using classical algebraic mechanisms, not an automorphic existence theorem. Written proofs, self-review, counterexamples, and exact finite regression checks accompany the constructions. No Lean certification, new foundational axiom, historical novelty, or Langlands reciprocity is claimed.
\end{abstract}
\newpage

\tableofcontents
\newpage

\section{Purpose and status}

The goal is to formalize the idea that a mathematical theory may contain two distinct ambient directions:
\[
\boxed{\text{transformation: change of representation}}
\qquad\text{and}\qquad
\boxed{\text{resolution: change of distinguishability}.}
\]
The distinction is operational rather than rhetorical. A Fourier transform, a duality, or a change of coordinates may alter representation while preserving the level of information being resolved. Coarsening a filtration, passing from modulus $p^{n+1}$ to modulus $p^n$, truncating a spectral expansion, or forgetting fine ramification data changes the level at which distinctions are visible.

The central hypothesis of this note is that a large class of mathematical constructions can be organized by a two-axis calculus in which:
\begin{enumerate}
    \item objects occur at resolutions;
    \item each resolution has its own transformation category;
    \item resolution changes induce functors between those categories;
    \item mixed transformation-resolution squares can commute exactly, commute up to equivalence, or leave a structured residual;
    \item residuals can themselves be promoted to new seeds;
    \item compatible local or resolved data can sometimes be reconstructed into global objects.
\end{enumerate}

\paragraph{Status convention.}
Results tagged ``proved here'' have written proofs under the stated hypotheses. They are not asserted to be Lean-checked. Results tagged ``standard machinery'' are direct applications of established category, order, or rewriting theory. Statements tagged ``program theorem'' are precise targets whose Komorebi-specific hypotheses remain to be supplied in applications. No claim of mathematical novelty is made merely because standard machinery is expressed in Komorebi language.

\paragraph{Categorical regimes and size.}
The initial indexed construction uses small ordinary categories and a strict contravariant functor. Every section involving functorial cofibers, higher coherence, stable envelopes, or stable Verdier quotients uses small stable $\infty$-categories, or a specified enhancement presenting them. All subcategories in closure constructions are full and replete. ``Stable'' means closed under finite limits and colimits, while ``thick'' additionally means closed under retracts. A stable seed zoo is not automatically idempotent-complete. Thick results are applied to a separately designated thick zoo or to its idempotent completion. Homotopy categories support the corresponding triangulated statements, not arbitrary canonical choices of cones. A fixed small presentation or larger ambient universe is understood whenever object collections are used as sets.

\paragraph{What the two axes do not assert.}
An admissible transformation need not be invertible or information-preserving. A coarsening functor need not have an inverse. A factorization in a constructed fibration does not establish a canonical resolution system for every mathematical theory. These are separate semantic requirements, not consequences of the notation.

\section{Resolution systems}

\subsection{The resolution base}

\begin{definition}[Resolution category]
A \emph{resolution category} is a small poset $(R,\le)$ regarded as a thin category $\Rcat$. An arrow $r\to s$ exists exactly when $r\le s$. We interpret $s$ as at least as fine as $r$.
\end{definition}

The poset assumption is intentionally restrictive for the first formalization. A later version may replace $\Rcat$ by a general category, a site, a filtered category, or a multi-parameter resolution geometry.

\begin{definition}[Indexed transformation system]
An \emph{indexed transformation system} over $\Rcat$ is a strict functor
\[
F:\Rcat^{op}\longrightarrow\Cat.
\]
For $r\in R$, write
\[
\Tcat_r:=F(r).
\]
For $r\le s$, write
\[
C_{s\to r}:=F(r\le s):\Tcat_s\longrightarrow\Tcat_r.
\]
We call $C_{s\to r}$ the \emph{coarsening functor}.
\end{definition}

Strictness gives
\[
C_{r\to r}=\id_{\Tcat_r}
\]
and, for $r\le s\le t$,
\[
C_{t\to r}=C_{s\to r}\circ C_{t\to s}.
\]
A pseudofunctor version is more natural in some applications. The strict version isolates the main structure without coherence bookkeeping.

\begin{definition}[Resolved object]
A \emph{resolved object} is a pair $(r,X)$ with $X\in\Ob(\Tcat_r)$.
\end{definition}

\subsection{Resolution equivalence}

Let $X,Y\in\Tcat_t$ and let $r\le t$.

\begin{definition}[Indistinguishability at resolution $r$]
We write
\[
X\sim_r Y
\]
when
\[
C_{t\to r}(X)\cong C_{t\to r}(Y)
\]
in $\Tcat_r$.
\end{definition}

For a fixed ambient fine level $t$, each $\sim_r$ is an equivalence relation on isomorphism classes of objects of $\Tcat_t$.

\begin{theorem}[Resolution monotonicity]\label{thm:resolution-monotonicity}
Let $r\le s\le t$ and $X,Y\in\Tcat_t$. If
\[
X\sim_sY,
\]
then
\[
X\sim_rY.
\]
\end{theorem}

\begin{proof}
An isomorphism
\[
C_{t\to s}(X)\cong C_{t\to s}(Y)
\]
is carried by the functor $C_{s\to r}$ to an isomorphism
\[
C_{s\to r}C_{t\to s}(X)
\cong
C_{s\to r}C_{t\to s}(Y).
\]
By functorial coherence this is
\[
C_{t\to r}(X)\cong C_{t\to r}(Y).
\]
Hence $X\sim_rY$.
\end{proof}

\begin{corollary}[Nested distinguishability]\label{cor:nested}
For $r\le s\le t$, the equivalence relation $\sim_s$ refines $\sim_r$.
\end{corollary}

This is the first exact form of the statement that coarsening may erase distinctions but cannot create distinctions that were absent at finer resolution.

\subsection{Refinement as a fiber rather than a function}

\begin{definition}[Lift set]
For $r\le s$ and $X\in\Tcat_r$, define the set of isomorphism classes of lifts
\[
\Lift_{r\to s}(X)
:=
\left\{
Y\in\Ob(\Tcat_s):C_{s\to r}(Y)\cong X
\right\}/\cong.
\]
\end{definition}

Three basic cases occur:
\[
|\Lift_{r\to s}(X)|=0,
\qquad
|\Lift_{r\to s}(X)|=1,
\qquad
|\Lift_{r\to s}(X)|>1.
\]
They correspond respectively to obstruction to refinement, uniqueness of the fine object's isomorphism class, and hidden fine-scale multiplicity. The middle case does not imply a contractible space of refinements. It forgets the identification $C_{s\to r}Y\cong X$ and all automorphisms. When these matter, use the groupoid of pairs $(Y,\theta)$ with $\theta:C_{s\to r}Y\xrightarrow{\sim}X$, and isomorphisms commuting with $\theta$. In the enhanced setting this is the homotopy fiber of the functor on cores.

\begin{definition}[Resolution branching number]
When $\Lift_{r\to s}(X)$ is finite, define
\[
b_{r\to s}(X):=|\Lift_{r\to s}(X)|.
\]
When $b_{r\to s}(X)>0$, its logarithm
\[
H_{r\to s}(X):=\log b_{r\to s}(X)
\]
may be regarded as a crude \emph{resolution information} measuring how much ambiguity is hidden by coarsening.
\end{definition}

The empty lift set is recorded as an obstruction, not as an ordinary finite entropy. This quantity is only a first model. In probabilistic, measurable, or derived settings one would replace cardinality by entropy, measure, dimension, rank, or a homotopical invariant.

\section{The total two-axis category}

\subsection{Grothendieck construction}

\begin{construction}[Total Komorebi category]\label{cons:groth}
Given $F:\Rcat^{op}\to\Cat$, define the category
\[
\K:=\int_{\Rcat}F
\]
as follows.

Its objects are pairs $(r,X)$ with $X\in\Tcat_r$.

For $r\le s$, a morphism
\[
(r,X)\longrightarrow(s,Y)
\]
is a morphism
\[
\alpha:X\longrightarrow C_{s\to r}(Y)
\]
in $\Tcat_r$. If $r\nleq s$, there is no such morphism in this directed version.
\end{construction}

Given
\[
(r,X)\xrightarrow{\alpha}(s,Y)
\quad\text{and}\quad
(s,Y)\xrightarrow{\beta}(t,Z),
\]
define
\[
\beta\circ_{\K}\alpha
:=
C_{s\to r}(\beta)\circ\alpha:
X\to C_{t\to r}(Z).
\]

\begin{theorem}[Total category]\label{thm:total-category}
The construction above defines a category $\K$.
\end{theorem}

\begin{proof}
Identity at $(r,X)$ is $\id_X:X\to C_{r\to r}(X)=X$. For associativity, let
\[
(r,X)\xrightarrow{\alpha}(s,Y)
\xrightarrow{\beta}(t,Z)
\xrightarrow{\gamma}(u,W).
\]
One composite gives
\[
C_{t\to r}(\gamma)\circ C_{s\to r}(\beta)\circ\alpha.
\]
The other gives the same expression by functoriality and
\[
C_{t\to r}=C_{s\to r}C_{t\to s}.
\]
Thus composition is associative.
\end{proof}

Define
\[
\pi:\K\to\Rcat,
\qquad
\pi(r,X)=r.
\]
The fiber over $r$ is canonically equivalent to $\Tcat_r$.

\subsection{Canonical factorization}

For $r\le s$ and $Y\in\Tcat_s$, define the canonical morphism
\[
\kappa_{r,s,Y}:
(r,C_{s\to r}Y)\longrightarrow(s,Y)
\]
represented by the identity map
\[
\id_{C_{s\to r}Y}.
\]

\begin{theorem}[Two-axis factorization]\label{thm:factorization}
Every morphism
\[
f:(r,X)\longrightarrow(s,Y)
\]
in $\K$ factors uniquely as
\[
(r,X)
\xrightarrow{f_v}
(r,C_{s\to r}Y)
\xrightarrow{\kappa_{r,s,Y}}
(s,Y),
\]
where $f_v$ lies entirely in the fiber $\Tcat_r$.
\end{theorem}

\begin{proof}
By definition, $f$ is exactly a morphism
\[
\alpha:X\to C_{s\to r}Y
\]
in $\Tcat_r$. Take $f_v=\alpha$. Composition with $\kappa_{r,s,Y}$ is represented by
\[
\id_{C_{s\to r}Y}\circ\alpha=\alpha,
\]
so the composite is $f$. Any factorization through the fixed cartesian target $(r,C_{s\to r}Y)$ must use the same fiber morphism $\alpha$, hence is unique.
\end{proof}

\begin{remark}
This is the formal version of the two-axis statement
\[
\boxed{
\text{mixed movement}
=
\text{within-resolution transformation}
+
\text{resolution movement}.
}
\]
Categorically, $\pi:\K\to\Rcat$ is the Grothendieck fibration associated with $F$. Indeed, a map $(u,Z)\to(s,Y)$ above $u\le r\le s$ is exactly a map $Z\to C_{s\to u}Y=C_{r\to u}C_{s\to r}Y$. This uniquely factors through $\kappa_{r,s,Y}$ as a morphism to $(r,C_{s\to r}Y)$ above $u\le r$, which is the cartesian lifting property. Notice the direction: the lift starts from an already specified fine object $Y$. It is not a procedure that selects a refinement of an arbitrary coarse object.
\end{remark}

\section{Resolution geometry}

The nested equivalence relations from \cref{cor:nested} induce a genuine geometry when the resolution set is linearly ordered.

\begin{definition}[Separation depth]
Let $S$ be a set of states, with equivalence relations $\sim_n$ for every $n\in\mathbb N=\{0,1,\ldots\}$. Assume
\[
\sim_{n+1}\ \subseteq\ \sim_n,
\qquad \sim_0=S\times S.
\]
Define
\[
\nu(X,Y):=\sup\{n\in\mathbb N:X\sim_nY\}\in\mathbb N\cup\{\infty\}.
\]
The universal zeroth level makes this set nonempty. Assume separation, meaning that $X\sim_nY$ for every $n$ implies $X=Y$. The states can, for example, be coherent towers equipped with the chosen levelwise equivalence relations. No finiteness of the individual resolutions is required.
\end{definition}

\begin{theorem}[Resolution ultrametric]\label{thm:ultrametric}
Define
\[
d_R(X,Y):=
\begin{cases}
0,&X=Y,\\
2^{-\nu(X,Y)},&X\ne Y.
\end{cases}
\]
Then $d_R$ is an ultrametric:
\[
d_R(X,Z)\le\max\{d_R(X,Y),d_R(Y,Z)\}.
\]
\end{theorem}

\begin{proof}
Let
\[
m=\min\{\nu(X,Y),\nu(Y,Z)\}.
\]
For every $n\le m$, we have $X\sim_nY$ and $Y\sim_nZ$. Since $\sim_n$ is an equivalence relation, $X\sim_nZ$. Hence
\[
\nu(X,Z)\ge m.
\]
Therefore
\[
2^{-\nu(X,Z)}
\le
2^{-m}
=
\max\{2^{-\nu(X,Y)},2^{-\nu(Y,Z)}\}.
\]
The identity and symmetry axioms are immediate, and separation gives positivity for distinct states. Without separation the same formula, with $2^{-\infty}=0$, is an ultrapseudometric.
\end{proof}

Thus the resolution axis can carry an intrinsic non-Archimedean geometry. Two objects are close when many levels of resolution are required before they can be distinguished.

\begin{definition}[First distinguishing scale]
When finite, define
\[
\delta(X,Y):=\nu(X,Y)+1.
\]
This is the first resolution at which $X$ and $Y$ become distinguishable.
\end{definition}

This quantity may be useful in $p$-adic, dyadic, symbolic, filtration, and machine-representation settings.

\section{Seeds and generated zoos}

\subsection{Seed signatures}

Let $U$ be a set of isomorphism classes of resolved objects under consideration. Smallness is imposed only to use ordinary powerset fixed-point arguments.

\begin{definition}[Seed signature]
A \emph{seed signature} on $U$ is data
\[
\mathfrak S=(\Seeds_0,\mathscr T,\mathscr C,\mathscr L,\mathscr D),
\]
where:
\begin{itemize}
    \item $\Seeds_0\subseteq U$ is the initial seed set;
    \item $\mathscr T$ is a collection of partial unary transformation constructors on $U$;
    \item $\mathscr C$ is the collection of allowed coarsening constructors;
    \item $\mathscr L$ assigns to an object its admissible refinement classes;
    \item $\mathscr D$ is a collection of partial residual constructors.
\end{itemize}
\end{definition}

The use of partial constructors is intentional. An operation need not be defined on every object or every resolution.

\begin{definition}[One-step zoo operator]
For $A\subseteq U$, define
\[
\Phi(A):=
A
\cup \mathscr T[A]
\cup \mathscr C[A]
\cup \mathscr L[A]
\cup \mathscr D[A],
\]
where each term denotes all defined outputs of the corresponding constructors applied to inputs from $A$.
\end{definition}

\begin{lemma}[Monotonicity]\label{lem:phi-monotone}
If $A\subseteq B$, then
\[
\Phi(A)\subseteq\Phi(B).
\]
\end{lemma}

\begin{proof}
Every constructor output available from inputs in $A$ is also available from the same inputs viewed as elements of $B$.
\end{proof}

\begin{theorem}[Seed zoo existence]\label{thm:zoo-existence}
There exists a least $\Phi$-closed set containing $\Seeds_0$. Equivalently, there is a least fixed point of the closure process above $\Seeds_0$.
\end{theorem}

\begin{proof}
The powerset $\mathcal P(U)$ is a complete lattice. The map
\[
\Psi(A):=\Seeds_0\cup\Phi(A)
\]
is monotone by \cref{lem:phi-monotone}. By the Knaster--Tarski fixed-point theorem, $\Psi$ has a least fixed point. This is exactly the least constructor-closed set containing $\Seeds_0$.
\end{proof}

\begin{definition}[Generated zoo]
The least fixed point in \cref{thm:zoo-existence} is denoted
\[
\Zoo(\Seeds_0).
\]
\end{definition}

\begin{proposition}[Finitary iteration]\label{prop:finitary}
Assume each constructor is given by a fixed input/output relation with finitely many inputs, independent of the stage $A$. Then $\Phi$ preserves increasing unions of $\omega$-chains. Then
\[
\Zoo(\Seeds_0)
=
\bigcup_{n<\omega}\Psi^n(\varnothing).
\]
\end{proposition}

\begin{proof}
For a constructor output from $\bigcup_n A_n$, its finitely many inputs all belong to a common $A_N$. The output therefore already belongs to $\Phi(A_N)$. This proves preservation of increasing unions. The union on the right is a fixed point because $\Psi$ preserves that union. It contains $\Seeds_0$ and is contained in every fixed point by induction on $n$, hence is least.
\end{proof}

This turns the seed calculus into an explicit generative process:
\[
\Seeds_0
\subseteq
\Seeds_1
\subseteq
\Seeds_2
\subseteq\cdots,
\qquad
\Seeds_{n+1}=\Psi(\Seeds_n).
\]

\subsection{Zoo growth}

\begin{definition}[Growth function]
When each finite stage contains finitely many isomorphism classes, define
\[
g_{\Seeds_0}(n):=|\Psi^n(\varnothing)|.
\]
\end{definition}

The asymptotic class of $g(n)$ may distinguish finite, polynomial, exponential, and faster-growing presentations. It depends on the chosen generator signature and cost convention, not merely on the generated category. This is a proposed invariant of a generated mathematical grammar rather than a theorem of the present note.

\section{Transformation-resolution exchange}

\subsection{Resolution-indexed transformations}

\begin{definition}[Lax resolution-compatible transformation]\label{def:laxtrans}
A \emph{lax resolution-compatible transformation} consists of endofunctors
\[
A_r:\Tcat_r\to\Tcat_r
\]
for every $r\in R$, together with natural transformations
\[
\eta^A_{s,r}:
A_rC_{s\to r}
\Longrightarrow
C_{s\to r}A_s
\qquad(r\le s),
\]
satisfying identity and composition coherence:
\[
\eta^A_{r,r}=\id
\]
and, for $r\le s\le t$,
\[
\eta^A_{t,r}
=
\bigl(C_{s\to r}\eta^A_{t,s}\bigr)
\circ
\bigl(\eta^A_{s,r}C_{t\to s}\bigr).
\]
If every $\eta^A_{s,r}$ is a natural isomorphism, $A$ is called \emph{invertibly resolution-compatible}. This is distinct from exactness of a functor between stable categories.
\end{definition}

The square is
\[
\begin{tikzcd}[column sep=large,row sep=large]
A_rC_{s\to r}(X) \arrow[r,"\eta^A_{s,r,X}"] & C_{s\to r}A_s(X).
\end{tikzcd}
\]
It compares transforming after coarsening with coarsening after transforming.

\begin{theorem}[Lift to the total category]\label{thm:lift-total}
Every lax resolution-compatible transformation $A$ induces an endofunctor
\[
\widehat A:\K\to\K
\]
defined on objects by
\[
\widehat A(r,X)=(r,A_rX).
\]
\end{theorem}

\begin{proof}
A morphism $(r,X)\to(s,Y)$ in $\K$ is represented by
\[
\alpha:X\to C_{s\to r}Y.
\]
Define $\widehat A(\alpha)$ to be
\[
A_rX
\xrightarrow{A_r(\alpha)}
A_rC_{s\to r}Y
\xrightarrow{\eta^A_{s,r,Y}}
C_{s\to r}A_sY.
\]
Identity preservation follows from $\eta^A_{r,r}=\id$. Composition preservation is exactly the coherence condition in \cref{def:laxtrans}, together with naturality of $\eta^A$.
\end{proof}

Thus coherent exchange data suffice to assemble the family into a functor of the total system. A converse requires the functor to lie over the resolution base and to have the specified fiberwise restrictions.

\subsection{Mixed residuals}

\begin{definition}[Stable mixed residual]
Assume each $\Tcat_r$ is a small stable $\infty$-category and the coarsening and transformation functors are exact, with homotopy-coherent exchange data. For $r\le s$ and $X\in\Tcat_s$, define
\[
\Delta^A_{s,r}(X)
:=
\Cofib\bigl(\eta^A_{s,r,X}\bigr)
\in\Tcat_r.
\]
\end{definition}

Then there is a cofiber sequence
\[
A_rC_{s\to r}X
\longrightarrow
C_{s\to r}A_sX
\longrightarrow
\Delta^A_{s,r}(X).
\]

\begin{proposition}[Residual vanishing criterion]\label{prop:residual-vanish}
In a stable fiber,
\[
\Delta^A_{s,r}(X)\simeq0
\]
if and only if
\[
\eta^A_{s,r,X}
\]
is an equivalence.
\end{proposition}

\begin{proof}
In a stable category a morphism is an equivalence exactly when its cofiber is a zero object.
\end{proof}

In an abelian fiber one may instead record
\[
\Delta^-_{s,r}(X)=\ker\eta^A_{s,r,X},
\qquad
\Delta^+_{s,r}(X)=\operatorname{coker}\eta^A_{s,r,X}.
\]
These vanish simultaneously exactly when the comparison is an isomorphism.

\begin{definition}[Residual promotion]
If $X$ is a seed-generated object and $\Delta^A_{s,r}(X)$ is nonzero, the operation
\[
X\longmapsto\Delta^A_{s,r}(X)
\]
is called \emph{residual promotion} when the residual is added to the active seed set.
\end{definition}

This creates the recursive pattern
\[
\text{seed}
\to
\text{mixed square}
\to
\text{residual}
\to
\text{new seed}
\to
\text{expanded zoo}.
\]

\section{Normal forms for mixed words}

To make the two-axis calculus computational, mixed transformation-resolution expressions should admit controlled rewriting.

\subsection{Typed word system}

Consider a typed alphabet containing transformation symbols $T$ and resolution symbols $C$. A mixed word is a composable expression such as
\[
W=T_1C_1T_2T_3C_2T_4.
\]
Assume every well-typed adjacent transformation/coarsening inversion admits at least one length-preserving exchange rule of the form
\[
TC\;\Rightarrow\;C\,T',
\]
where $T'$ is the correctly reindexed transformation at the new resolution.

\begin{definition}[Inversion number]
For a word $W$, let $\iota(W)$ be the number of ordered pairs $(T,C)$ in which a transformation symbol occurs to the left of a resolution symbol when the chosen normal form requires all resolution symbols first.
\end{definition}

\begin{lemma}[Exchange termination]\label{lem:termination}
Every exchange rewrite
\[
TC\Rightarrow CT'
\]
reduces $\iota$ by one. Therefore every sequence of such rewrites terminates.
\end{lemma}

\begin{proof}
The rewritten adjacent pair contributes one inversion before the rewrite and none afterward. The relative order of every other transformation-resolution pair is unchanged.
\end{proof}

\begin{theorem}[Exact two-axis normal form]\label{thm:normal-form}
Assume the exchange rewriting system is locally confluent. Then every well-typed mixed word has a unique normal form up to the equality notion of the rewriting system. The normal form has all resolution operations grouped before all within-resolution transformations.
\end{theorem}

\begin{proof}
Induct on the inversion number. Every rewrite lowers it, so normal forms are reached. If two paths start with different rewrites, local confluence gives joining paths from their first successors to a common word. Normalize that word. The induction hypothesis at each first successor identifies its normal form with the common descendant's normal form. Thus the two original paths have the same normal form. Since every adjacent inversion has a rule, a normal word has no inversion and therefore has all coarsening symbols to the left of all transformation symbols. This is the finite-height form of Newman's argument.
\end{proof}

\status{proved here from standard rewriting theory under local confluence}

The nontrivial application problem is to prove local confluence for the concrete labelled signature. This theorem concerns word syntax, not equality of the semantic comparison maps along different paths. The pure two-sort grammar is treated more explicitly in \cref{thm:event-ideals}.

\subsection{Split residual expansion}

Exact exchange is too restrictive for many applications. A first rigorous residual model can be obtained in an additive setting.

\begin{definition}[Split residual exchange]
A chosen exchange morphism $\eta:TC\Rightarrow CT'$ is \emph{split residual} if it is naturally identified with the projection in a decomposition
\[
TC(X)\cong CT'(X)\oplus E_{T,C}(X).
\]
Here $E_{T,C}$ is the split \emph{source-side remainder}. It is not denoted by $\Delta$, because in the stable convention of this paper
\[
\Cofib(\eta_X)\simeq E_{T,C}(X)[1],
\qquad E_{T,C}(X)\simeq\Delta_{T,C}(X)[-1].
\]
\end{definition}

\begin{theorem}[Finite split residual expansion]\label{thm:split-expansion}
Fix a finite principal normalization path for a well-typed mixed word. Assume every step has a chosen split residual exchange and every surrounding operator preserves finite direct sums. Then
\[
W(s)\cong W_{\mathrm{nf}}(s)\oplus\bigoplus_{i=1}^m E_i(s),
\]
where $E_i(s)$ is the surrounding context applied to the source-side remainder of step $i$. In the stable case, $E_i(s)\simeq D_i[-1]$ for the step cofibers $D_i$ defined below.
\end{theorem}

\begin{proof}
At step $i$, the split exchange and additivity of the surrounding context give
$X_{i-1}\cong X_i\oplus E_i(s)$. Substituting these decompositions successively yields
$X_0\cong X_m\oplus E_m(s)\oplus\cdots\oplus E_1(s)$.
The direct sum is finite because the chosen principal path is finite. The identification with shifted cofibers follows from the cofiber of a split projection.
\end{proof}

This proof does not recursively normalize the residual branches. Such a recursive residual forest needs its own well-founded measure. Termination of the principal word alone does not prove termination of every newly generated residual word.

\begin{definition}[Normalization path]
Let $\mathcal S$ be a stable $\infty$-category, or a stable enhancement with functorial cofibers. A \emph{normalization path} for a mixed word $W$ evaluated on $s$ is a finite chain
\[
X_0=W(s)\xrightarrow{f_1}X_1\xrightarrow{f_2}\cdots\xrightarrow{f_m}X_m=W_{\mathrm{nf}}(s),
\]
where every $f_i$ is obtained by applying an exact surrounding context to one elementary exchange morphism. Write
\[
D_i:=\Cofib(f_i),
\qquad
D_{\mathrm{tot}}:=\Cofib(X_0\to X_m).
\]
\end{definition}

\begin{theorem}[Stable residual filtration along a normalization path]\label{thm:stable-filtration}
For every finite normalization path in a stable enhancement there is a canonical finite filtration
\[
0=F_0\longrightarrow F_1\longrightarrow\cdots\longrightarrow F_m=D_{\mathrm{tot}}
\]
with cofiber sequences
\[
F_{i-1}\longrightarrow F_i\longrightarrow D_i.
\]
Consequently the associated graded pieces of the total normalization defect are precisely the transported elementary exchange residuals $D_i$.
\end{theorem}

\begin{proof}
Define
\[
F_i:=\Cofib(X_0\longrightarrow X_i).
\]
Then $F_0\simeq0$ and $F_m=D_{\mathrm{tot}}$. For the composable maps
\[
X_0\longrightarrow X_{i-1}\xrightarrow{f_i}X_i,
\]
the octahedral cofiber sequence gives
\[
\Cofib(X_0\to X_{i-1})
\longrightarrow
\Cofib(X_0\to X_i)
\longrightarrow
\Cofib(f_i).
\]
Thus
\[
F_{i-1}\longrightarrow F_i\longrightarrow D_i
\]
is a cofiber sequence. Since $f_i$ is obtained by placing an elementary exchange morphism inside an exact context, exactness identifies $D_i$ with the same context applied to the elementary exchange cofiber. Iterating gives the stated filtration.
\end{proof}

\status{proved for a chosen normalization path in a stable enhancement}

\begin{corollary}[Source-to-normal-form residual control]\label{cor:source-normal-control}
There is a cofiber sequence
\[
D_{\mathrm{tot}}[-1]\longrightarrow W(s)\longrightarrow W_{\mathrm{nf}}(s),
\]
and $D_{\mathrm{tot}}[-1]$ carries the shifted filtration with successive pieces $D_i[-1]$. Hence the failure of the original word to equal its normal form is an iterated extension of shifted elementary exchange residuals.
\end{corollary}

The theorem is path-relative. If two normalization paths are related by coherent higher exchange cells, their total comparison maps may agree up to homotopy, but path-independence of the induced filtration is a separate problem. In an ordinary triangulated category without an enhancement, the octahedral axiom gives existence of such filtrations but not a canonical functorial choice.

\begin{problem}[Order-independent residual filtration]\label{prob:path-independent}
Identify an intrinsic replacement for the family of path-relative filtrations. Mere higher coherence cannot force their ordered stages to be equivalent: \cref{prop:filtration-no-go} gives a counterexample. The residual diagram of \cref{thm:ideal-residual} is a replacement when a coherent event diagram is supplied. Additional hypotheses for equivalence of selected ordered filtrations must constrain the intermediate objects, not only the comparison maps.
\end{problem}

\section{Normalization histories and intrinsic residual diagrams}\label{sec:histories}

There are three different invariance questions. A word may have a unique normal-form \emph{object}, the normalization paths may induce the same \emph{comparison map}, and the cofiber of that map may carry several different \emph{ordered filtrations}. None of these statements should be silently substituted for the next.

\subsection{A no-go theorem for ordered filtration independence}

\begin{proposition}[Coherent diamonds do not determine one ordered filtration]\label{prop:filtration-no-go}
Even a strictly commuting bicartesian square of objects in a stable category can yield two normalization paths with the same total comparison map and inequivalent ordered residual filtrations.
\end{proposition}

\begin{proof}
Work in the stable category of perfect complexes over a field $k$. Put $M=k[0]$ and $N=k^2[0]$, and consider the square
\[
\begin{tikzcd}[column sep=large,row sep=large]
0 \arrow[r] \arrow[d] & M \arrow[d,hook] \\
N \arrow[r,hook] & M\oplus N.
\end{tikzcd}
\]
The maps into $M\oplus N$ are the summand inclusions. This square is a pushout, hence also a pullback. Both composites are the unique map $0\to M\oplus N$.

The residual filtrations of its two boundary paths are
\[
0\longrightarrow M\longrightarrow M\oplus N,
\qquad
0\longrightarrow N\longrightarrow M\oplus N.
\]
An equivalence of these objects of $\operatorname{Fun}([2],\operatorname{Perf}(k))$ would induce an equivalence $M\simeq N$ at the middle index. That is impossible because their degree-zero cohomology dimensions are $1$ and $2$. Nevertheless both total residuals are $M\oplus N$.
\end{proof}

\begin{example}[The counterexample occurs inside the two-axis grammar]\label{ex:typed-diamond}
The preceding square is not merely an unrelated diagram. Let the resolution base be $0<1<2$, put
\[
\Tcat_0=\Tcat_1=\Tcat_2=\operatorname{Perf}(k)\times\operatorname{Perf}(k),
\]
and take every coarsening functor to be the identity between labelled copies. Define the exact coordinate projections
\[
P_1(x,y)=(x,0),\qquad P_2(x,y)=(0,y).
\]
Use transformation families
\[
A_0=P_2,\quad A_1=A_2=\id,
\qquad
B_0=B_1=P_1,\quad B_2=\id.
\]
For $A$, the comparison across $1\to0$ is the inclusion $P_2\Rightarrow\id$, and across $2\to1$ it is the identity. For $B$, these choices are respectively the identity and $P_1\Rightarrow\id$. The long comparisons are their composites. Units and transitivity therefore hold.

On the seed $x=(k[0],k^2[0])$, the word
\[
W=A_0C_{1\to0}B_1C_{2\to1}
\]
evaluates to zero. Exchanging its left inversion first yields $P_1x$. Exchanging its right inversion first yields $P_2x$. The second exchange in either path yields $x$. Both paths then perform the same identity exchange to reach the normal word
\[
C_{1\to0}C_{2\to1}A_2B_2.
\]
Thus the complete paths have residual filtrations
\[
0\to P_1x\to x\xrightarrow{\id}x,
\qquad
0\to P_2x\to x\xrightarrow{\id}x.
\]
Their first nonzero stages are not equivalent. The transformations, contexts, and coarsenings are exact, and all comparison data are coherent. Consequently extra coherence alone cannot repair the general claim of ordered filtration independence.
\end{example}

\status{general ordered-filtration independence refuted; the path-relative theorem remains valid}

The correction is not to discard the filtration. It is to retain the collection of compatible intermediate stages in a diagram richer than a single linear order.

\subsection{What suffices for independence of the total residual}

The normal-form endpoint alone does not determine a residual. For example, the maps $\id_k$ and $0:k\to k$ have the same source and target, but cofibers $0$ and $k\oplus k[1]$. Comparison maps therefore belong to the data.

\begin{theorem}[Semantic diamond criterion]\label{thm:semantic-diamond}
Let a rewrite graph have a height $h$ in $\mathbb N$ that decreases on every edge, and suppose every vertex $x$ has a unique reachable normal vertex $n(x)$. Assign an object $X_x$ in a stable $\infty$-category to each vertex and a morphism to each edge. Assume every elementary peak
\[
x\xrightarrow{e}y,\qquad x\xrightarrow{f}z
\]
has joining paths $a:y\leadsto u$ and $b:z\leadsto u$ whose semantic composites satisfy
\[
[a]\,[e]=[b]\,[f]\quad\text{in the homotopy category}.
\]
Then every normalization path from $x$ to $n(x)$ induces the same homotopy class of comparison map. In particular, the total residuals of any two such paths are equivalent.
\end{theorem}

\begin{proof}
Induct on $h(x)$. A normal vertex has only the empty normalization path. Otherwise take two normalization paths with initial edges $e$ and $f$. Choose the joining paths in the hypothesis and any normalization path $t:u\leadsto n(x)$. Such a path exists, and its terminal vertex is $n(x)$ by uniqueness of normal forms.

Since $h(y),h(z)<h(x)$, the induction hypothesis replaces the two original suffixes by $ta$ and $tb$ respectively without changing their homotopy classes. The full comparison classes are therefore
\[
[t][a][e]=[t][b][f].
\]
This proves equality for all normalization paths from $x$. Equality in the homotopy category gives a homotopy between the corresponding maps with fixed endpoints. A homotopy identifies the two arrows up to equivalence, and the cofiber functor preserves that equivalence.
\end{proof}

\begin{remark}[Equivalence is not a chosen canonical identification]
The theorem asserts existence of equivalent total cofibers. A specified equivalence requires a specified homotopy between comparison maps. Independence of all such choices requires higher coherence. Neither conclusion identifies the intermediate ordered filtrations in \cref{prop:filtration-no-go}.
\end{remark}

\subsection{The event poset of pure exchange}

The pure exchange grammar has enough combinatorial structure to index all normalization histories without choosing their order. This subsection concerns only adjacent swaps $TC\Rightarrow CT$ that retain the identity of each symbol. More elaborate rules that duplicate, delete, or replace generators are not included.

\begin{definition}[Exchange event poset]\label{def:event-poset}
Fix a finite word $W$ in two types of labelled symbols. Number its transformation symbols $T_1,\ldots,T_a$ and its coarsening symbols $C_1,\ldots,C_b$ in their original left-to-right order. An \emph{exchange event} is a pair
\[
(i,j)\in E(W)
\quad\Longleftrightarrow\quad
T_i\text{ occurs to the left of }C_j\text{ in }W.
\]
The event swaps the relative order of this pair exactly once. Define $P_W$ as the partial order generated by
\begin{align*}
(i,k)&\prec(i,j) &&\text{when }k<j\text{ and both events exist},\\
(h,j)&\prec(i,j) &&\text{when }h>i\text{ and both events exist}.
\end{align*}
These relations are acyclic: along a generating relation the transformation index decreases or the coarsening index increases, with at least one strict change. Let $J(P_W)$ be the poset of order ideals, ordered by inclusion.
\end{definition}

The first relation says that one transformation crosses the coarsenings from left to right. The second says that one coarsening crosses the transformations from right to left.

\begin{theorem}[Normalization states are order ideals]\label{thm:event-ideals}
For the pure adjacent-swap grammar, the following correspondences are bijective:
\begin{enumerate}
\item words reachable from $W$ and ideals $I\in J(P_W)$;
\item enabled elementary swaps at a reachable word and minimal elements of $P_W\setminus I$;
\item complete normalization histories and linear extensions of $P_W$.
\end{enumerate}
The ideal attached to a reachable word is the set of events already performed. The original word corresponds to $\varnothing$, and the normal word to $P_W$.
\end{theorem}

\begin{proof}
Every swap changes the order of exactly one transformation/coarsening pair, never to be reversed. The relative orders of the transformations among themselves and of the coarsenings among themselves remain fixed.

Before event $(i,j)$ can occur, every initially intervening $C_k$ with $k<j$ must already have crossed $T_i$, and every initially intervening $T_h$ with $h>i$ must already have crossed $C_j$. Thus all predecessors of a performed event have already occurred. The performed events form an ideal.

That ideal determines the reachable word uniquely. It specifies whether each originally inverted pair has crossed, while pairs originally of the form $C_jT_i$ never change order. Together with the fixed orders inside each type, this determines every pairwise order of labelled symbols.

Conversely, suppose an ideal $I$ has been performed and $(i,j)$ is a minimal event outside it. In the current word $T_i$ still precedes $C_j$. If another symbol lies between them, it is either some $C_k$ with $k<j$ or some $T_h$ with $h>i$. In the first case $(i,k)$ is an unperformed predecessor. In the second case $(h,j)$ is an unperformed predecessor. The relevant pair was already inverted in the original word, since swaps never turn an original $CT$ pair into $TC$. Both cases contradict minimality. Hence $T_i$ and $C_j$ are adjacent and the event is enabled.

Any ideal can now be constructed by adding its events in an order compatible with $P_W$. This proves surjectivity and the statement about edges. A complete path adds each event exactly once, so it is exactly a linear extension. Conversely each linear extension prescribes enabled swaps and therefore a complete path.
\end{proof}

\begin{lemma}[All histories are related by independent interchanges]\label{lem:linear-extension-swaps}
Any two linear extensions of a finite poset are connected by swaps of adjacent incomparable elements. Consequently any two pure normalization histories are related by commuting independent-event diamonds.
\end{lemma}

\begin{proof}
Let $e$ be the first element of the second linear extension. It is minimal in the poset. In the first extension, every element appearing before $e$ is incomparable with it: it cannot precede $e$ in the poset by minimality, and it cannot follow $e$ in the poset while occurring before it in a linear extension. Move $e$ left across these incomparable elements. The two extensions now begin with the same element. Delete it and apply induction to the remaining finite poset.

In the exchange poset, incomparable events share neither a transformation symbol nor a coarsening symbol. When both are enabled their adjacent swaps are disjoint. Performing them in either order gives the same word, namely the diamond for adjoining both events to the current ideal.
\end{proof}

It follows that the category of rewrite paths modulo these diamond relations is the thin category $J(P_W)$. The conclusion concerns the pure syntax. In a typed application, labels after crossing must be defined by the specified reindexing rules, and their semantic exchange cells must satisfy the corresponding comparisons.

\subsection{A diagram that retains the residual histories}

\begin{definition}[Coherent exchange-event diagram]
For a finite event poset $P$, a coherent exchange-event diagram is a functor
\[
X:NJ(P)\longrightarrow\mathcal S
\]
into a stable $\infty$-category. Here $N$ denotes the nerve. Its value $X_I$ is the semantic state after the ideal $I$ of events has been completed. Values on cover arrows realize the elementary exchanges, and its higher simplices supply the compatibility data.
\end{definition}

Strictly commuting exchange diagrams in a concrete enhanced model, such as complexes, supply examples. Merely checking squares in the homotopy category does not automatically supply all the higher data in this definition.

\begin{theorem}[Ideal-indexed residual diagram]\label{thm:ideal-residual}
A coherent exchange-event diagram determines a functor
\[
\mathcal F_X:NJ(P)\longrightarrow\mathcal S,
\qquad
\mathcal F_X(I)=\Cofib(X_{\varnothing}\longrightarrow X_I),
\]
with the following properties:
\begin{enumerate}
\item $\mathcal F_X(\varnothing)\simeq0$;
\item along every maximal chain of ideals it restricts to the residual filtration of \cref{thm:stable-filtration};
\item for every cover $I\subset I\cup\{e\}$ there is a cofiber sequence
\[
\mathcal F_X(I)\longrightarrow\mathcal F_X(I\cup\{e\})
\longrightarrow\Cofib(X_I\to X_{I\cup\{e\}});
\]
\item its terminal object is the total residual $\Cofib(X_{\varnothing}\to X_P)$, independent of the linear extension used to reach it;
\item exact functors preserve the entire diagram, naturally in the specified coherent input.
\end{enumerate}
\end{theorem}

\begin{proof}
The arrows from the initial ideal $\varnothing$ induce a natural transformation from the constant diagram at $X_{\varnothing}$ to $X$. Take its cofiber pointwise in the stable functor category. This constructs $\mathcal F_X$ together with its coherence, not just an assignment of objects.

At $\varnothing$ the map is the identity, so its cofiber is zero. Restriction to a maximal chain gives exactly the pointwise formula defining the earlier path-relative filtration. For an inclusion of ideals, apply the cofiber sequence for the composable maps
\[
X_{\varnothing}\longrightarrow X_I\longrightarrow X_{I\cup\{e\}}.
\]
This gives the displayed cover cofiber sequence. At the terminal ideal the same formula gives the total residual. Finally an exact functor preserves these finite colimits pointwise and hence preserves the diagram and all its structure maps.
\end{proof}

\status{proved relative to a specified coherent exchange-event diagram}

This is the appropriate replacement for a single path-independent ordered filtration. In \cref{prop:filtration-no-go}, $P$ has two incomparable events, so $J(P)$ is the four-object diamond. Its two chains remember different intermediate residuals without disagreeing about the total residual. The richer diagram preserves information that the proposed identification of ordered filtrations would have erased.

\begin{corollary}[History-independent total obstruction]\label{cor:history-obstruction}
Let $q:\mathcal S\to\mathcal O$ be exact. For a fixed coherent exchange-event diagram, the object
\[
q\mathcal F_X(P)
\]
is independent of the normalization history. Along every maximal chain it has the filtration of \cref{thm:obstruction-localization}. Therefore its nonvanishing forces a nonzero elementary quotient residual on \emph{each} normalization history. The particular event witnessing nonvanishing may depend on the history.
\end{corollary}

\begin{proof}
The terminal object is fixed by \cref{thm:ideal-residual}. Apply \cref{thm:obstruction-localization} separately to the finite filtration supplied by each maximal chain.
\end{proof}

The converse is not asserted. Nonzero elementary residuals can participate in extensions whose total residual vanishes. Nor does the corollary exhibit one event that must witness the obstruction in every history.


\section{Constructing semantic exchange diagrams}\label{sec:strict-semantics}

The coherent diagram in \cref{thm:ideal-residual} was an input. We now construct it for a concrete enhanced realization of the typed pure grammar. This does not strictify arbitrary higher-categorical exchange data.

\begin{definition}[Strict chain-natural realization]\label{def:strict-chain-realization}
Fix a finite resolution chain $r_0\le\cdots\le r_b$. At each level choose a category $\mathcal E_j$ of complexes and a functor $L_j:\mathcal E_j\to\mathcal S_j$ into a stable $\infty$-category. A model example is bounded complexes of finitely generated projective modules mapping to their derived category. Specify chain-level functors
\[
 C_j:\mathcal E_j\to\mathcal E_{j-1},\qquad
 T_{i,j}:\mathcal E_j\to\mathcal E_j,
\]
which induce exact functors in the chosen stable realizations, and strict natural transformations
\[
 \chi_{i,j}:T_{i,j-1}C_j\Longrightarrow C_jT_{i,j}
 \qquad(1\le j\le b).
\]
The components of $\chi$ are actual chain maps and their naturality squares commute as chain maps. Functor composition and whiskering are taken in the ordinary strict 2-category of these models.
\end{definition}

Coarsening symbols occur in their fixed order $C_1,\ldots,C_b$. A transformation symbol with exactly $j$ coarsening symbols to its left is interpreted as $T_{i,j}$. Hence every reachable word represents a functor from $\mathcal E_b$ to $\mathcal E_0$. Moving $T_i$ across $C_j$ reindexes its level rather than creating a new untyped symbol. No rule in this subsection contracts coarsenings or introduces a new exchange relation.

\begin{theorem}[Strict semantic realization of the event diagram]\label{thm:strict-event-realization}
For a typed pure word $W$ in a strict chain-natural realization and a seed $s\in\mathcal E_b$, evaluation gives a functor
\[
 X_W:NJ(P_W)\longrightarrow\mathcal S_0.
\]
Its vertices evaluate the reachable words and its cover arrows evaluate the whiskered exchange maps. Thus the residual diagram in \cref{thm:ideal-residual} is constructed, rather than assumed, for this class of realizations.
\end{theorem}

\begin{proof}
First work in the ordinary category of functors $\mathcal E_b\to\mathcal E_0$ and natural transformations. Assign to a rewrite edge its exchange transformation, whiskered by the surrounding word. Two simultaneously enabled swaps are disjoint. The two composites around their diamond are equal by the interchange law for natural transformations. This equality uses strict naturality, not invertibility of $\chi$.

For ideals $I\subseteq J$, two paths adding the events in $J\setminus I$ are linear extensions of that interval's induced event order. The argument of \cref{lem:linear-extension-swaps} connects them by adjacent interchanges of incomparable events. Each interchange is one of the commuting diamonds just checked. Hence all such paths induce the same natural transformation. Evaluation therefore factors through the thin category $J(P_W)$.

Evaluate this ordinary functor on $s$, take its nerve, and compose with $L_0$. An ordinary functor gives all simplices of its nerve, so the output is a coherent diagram, not merely a diagram in a homotopy category. The prescribed vertices and edges are unchanged by this construction.
\end{proof}

\begin{remark}[Scope of coherence]
Strict models are sufficient, not necessary. This theorem does not infer strict naturality from a collection of commuting squares in a homotopy category. Arbitrary weak exchange cells, additional rewrite relations, and higher reindexing coherence still need their own realization theorem. For the fixed pure grammar, however, no independent braid relation is needed: the only changes of history exchange disjoint events.
\end{remark}

\section{Interaction residuals and event reconstruction}\label{sec:interaction-reconstruction}

An exchange square can commute coherently while changing the residual of one event after another event is performed. The appropriate measurement is a second cofiber. It measures failure of excision, not failure of commutativity.

\subsection{The interaction of two exchanges}

Let a specified coherent square in a stable $\infty$-category $\mathcal S$ be
\[
\begin{tikzcd}
 A \arrow[r,"u"] \arrow[d,"v"'] & B \arrow[d,"v'"]\\
 C \arrow[r,"u'"'] & D.
\end{tikzcd}
\]
Write $E=\Cofib(u)$ and $E'=\Cofib(u')$. The square induces $E\to E'$.

\begin{definition}[Interaction residual]\label{def:interaction-residual}
The interaction residual of this specified square is
\[
 \Theta=\Cofib\bigl(\Cofib(A\to B)\longrightarrow
                         \Cofib(C\to D)\bigr).
\]
The specified square, including its coherent homotopy, is part of the input.
\end{definition}

\begin{theorem}[Interaction is the total cofiber]\label{thm:interaction-cofiber}
There are natural equivalences
\[
 \Theta\simeq\Cofib\bigl(B\amalg_A C\longrightarrow D\bigr)
 \simeq\Cofib\bigl(\Cofib(A\to C)\longrightarrow
                         \Cofib(B\to D)\bigr).
\]
Consequently $\Theta\simeq0$ if and only if the square is a homotopy pushout, equivalently a homotopy pullback. An exact functor preserves these equivalences.
\end{theorem}

\begin{proof}
Cofibers are pushouts along maps to zero. Interchanging the finite colimits in the two successive pushout constructions identifies the first expression with the cofiber of the map from the punctured-square colimit $B\amalg_A C$ to $D$. Performing the same construction in the other coordinate gives the second expression. These identifications arise from the universal properties of the same finite diagrams and are natural in the given square.

A map in a stable category has zero cofiber exactly when it is an equivalence. Thus vanishing says precisely that the punctured-square colimit maps equivalently to $D$. Stability identifies pushout squares with pullback squares. Exact functors preserve zero objects and these finite colimits.
\end{proof}

In a chain model, interchanging coordinates includes the usual graded signs. The theorem is a natural equivalence, not an assertion that differently ordered cone matrices are literally identical. The cone conventions are recorded in \cite{StacksCones}; the finite-colimit argument supplies the enhanced statement used here.

\subsection{An event's residual before and after unrelated events}

Fix a finite poset $P$ and a coherent diagram $X:NJ(P)\to\mathcal S$. For $e\in P$, put
\[
 D_e=\{p\in P:p<e\},\qquad
 \downarrow e=D_e\cup\{e\}.
\]
The eligible contexts for $e$ form the interval of ideals
\[
 \mathcal A_e=\{I\in J(P):D_e\subseteq I,\ e\notin I\}.
\]
Its initial object is $D_e$. Define the residual functor on this interval by
\[
 E_e(I)=\Cofib(X_I\to X_{I\cup\{e\}}).
\]
For incomparable events $e,f$ enabled at $I$, write $\Theta_{e,f}(I)$ for the interaction residual of their elementary square.

\begin{theorem}[Criterion for context-independent event residuals]\label{thm:event-context-invariance}
The following conditions are equivalent.
\begin{enumerate}
\item Every elementary independent-event square of $X$ is a homotopy pushout.
\item Every $\Theta_{e,f}(I)$ vanishes.
\item For each event $e$, the natural comparison
\[
 \operatorname{const}_{\mathcal A_e}E_e(D_e)\longrightarrow E_e
\]
is an equivalence of diagrams.
\end{enumerate}
In this regime the successive cofibers of every normalization history are the event-labelled objects
\[
 d_e:=\Cofib(X_{D_e}\to X_{\downarrow e}),
\]
listed in the order in which that history performs the events. Their labels and equivalence classes are independent of context, but their ordering and extension data need not disappear.
\end{theorem}

\begin{proof}
The equivalence of (1) and (2) is \cref{thm:interaction-cofiber}. A cover in $\mathcal A_e$ adds an event $f$ incomparable with $e$. Indeed, predecessors of $e$ already lie in $I$, while a successor of $e$ cannot belong to an ideal omitting $e$. On this cover the map of event residuals has cofiber $\Theta_{e,f}(I)$. Condition (2) therefore makes every cover map an equivalence. Any inclusion in the finite interval factors into covers. The comparisons from its initial object consequently give the objectwise, hence diagrammatic, equivalence in (3).

Conversely, (3) makes every cover comparison of event residuals an equivalence, so its cofiber vanishes. The final statement follows by restricting the residual diagram to a maximal chain and applying the same comparison to each enabled event.
\end{proof}

\subsection{Reconstructing the full diagram from principal ideals}

The preceding criterion gives more than a common list of residuals. It makes the diagram recoverable from its values on the individual events together with their predecessors.

\begin{lemma}[Elementary excision implies ideal excision]\label{lem:ideal-excision}
If every elementary independent-event square of $X:NJ(P)\to\mathcal S$ is a homotopy pushout, then for all ideals $I,J$ the square
\[
\begin{tikzcd}
 X_{I\cap J}\arrow[r]\arrow[d] &X_I\arrow[d]\\
 X_J\arrow[r]&X_{I\cup J}
\end{tikzcd}
\]
is a homotopy pushout.
\end{lemma}

\begin{proof}
Elements of $I\setminus J$ and $J\setminus I$ are pairwise incomparable across the two sets. If $p\in I\setminus J$ satisfied $p\le q\in J\setminus I$, downward closure of $J$ would imply $p\in J$, a contradiction. The reverse comparison is excluded similarly.

Choose linear extensions for both exclusive sets, beginning from $I\cap J$. Adjoining their successive initial segments gives a rectangular grid of ideals. Each cell adjoins one enabled event from each side and is an elementary independent-event square. It is a pushout by hypothesis. Repeated pasting of pushout squares gives the outer square displayed above. Empty exclusive sets give identity squares and cause no exception.
\end{proof}

\begin{theorem}[Principal-ideal reconstruction]\label{thm:principal-reconstruction}
Let $P$ be finite and $\mathcal S$ stable. Write $\operatorname{Exc}_0(J(P),\mathcal S)$ for the full subcategory of coherent diagrams $F:NJ(P)\to\mathcal S$ such that $F(\varnothing)\simeq0$ and every elementary independent-event square is a homotopy pushout. Restriction along
\[
 j:P\longrightarrow J(P),\qquad e\longmapsto\downarrow e
\]
induces an equivalence
\[
 \operatorname{Exc}_0(J(P),\mathcal S)
 \simeq\operatorname{Fun}(NP,\mathcal S).
\]
Its inverse sends $G:NP\to\mathcal S$ to
\[
 (j_!G)(I)=\colim_{e\in I}G(e),
 \qquad (j_!G)(\varnothing)=0.
\]
All colimits are finite. In particular, for an excisive event diagram $X$ its reduced residual diagram satisfies
\[
 \boxed{\mathcal F_X(I)\simeq
 \colim_{e\in I}\mathcal F_X(\downarrow e).}
\]
\end{theorem}

\begin{proof}
The displayed formula is the pointwise left Kan extension along $j$. It follows directly from the colimit universal property: maps from $j_!G$ to $F$ are the same as compatible maps $G(e)\to F(\downarrow e)$. For a principal ideal $\downarrow e$, its indexing poset has terminal object $e$, so $(j_!G)(\downarrow e)\simeq G(e)$.

For arbitrary ideals $I,J$, every chain in $I\cup J$ lies wholly in $I$ or wholly in $J$. Its maximal element lies in one ideal and all its predecessors lie there too. Hence the nerve of their union is the union of their nerves along the intersection. Cocones on $G|_{I\cup J}$ are precisely compatible cocones on $G|_I$ and $G|_J$. The resulting colimit is the pushout of the two colimits over the colimit on $I\cap J$. Thus $j_!G$ is reduced and excisive.

Now let $F$ be reduced and excisive. We prove the counit $j_!j^*F\to F$ is an equivalence on an ideal $I$ by induction on $|I|$. The empty ideal is the zero object on both sides. Choose a maximal element $e\in I$. Put $I'=I\setminus\{e\}$. Then
\[
 I=I'\cup\downarrow e,\qquad I'\cap\downarrow e=D_e.
\]
The ideal-excision lemma expresses $F(I)$ as the pushout of $F(I')$ and $F(\downarrow e)$ over $F(D_e)$. The same holds for $j_!j^*F$. Induction handles $I'$ and $D_e$, and the terminal-object observation handles $\downarrow e$. Therefore the counit is an equivalence at $I$. The unit and counit prove the asserted equivalence of functor categories, including natural transformations.

Finally, the total-cofiber construction is exact on square diagrams. Applied to the cofiber sequence of diagrams $\operatorname{const}X_{\varnothing}\to X\to\mathcal F_X$, it shows that the elementary total cofibers of $X$ and $\mathcal F_X$ agree, since a constant square has zero total cofiber. Thus $\mathcal F_X$ is reduced and excisive, and the first part applies.
\end{proof}

This is finite-poset reconstruction by left Kan extension and excision. It is a theorem about the specified event diagram, not the missing arithmetic globalization theorem K15. Its general categorical ingredients are standard, and no historical novelty claim is made.

\begin{corollary}[Independent events split, dependencies retain extensions]\label{cor:independent-splitting}
Under \cref{thm:principal-reconstruction}, if $P$ is an antichain then
\[
 \mathcal F_X(P)\simeq\bigoplus_{e\in P}d_e.
\]
For a general event poset, a direct-sum formula does not follow.
\end{corollary}

\begin{proof}
For an antichain, each principal ideal is a singleton, $D_e=\varnothing$, and a finite colimit over the discrete poset is a finite direct sum. For the negative statement, take $P$ to be a two-element chain and a reduced diagram in $\operatorname{Perf}(\mathbb Z)$
\[
 0\longrightarrow\mathbb Z/p\xrightarrow{\,p\,}\mathbb Z/p^2
\]
for a prime $p$, where the second map sends the class of $x$ to the class of $px$. Both successive cofibers are $\mathbb Z/p$ in degree zero, but the terminal object is not $\mathbb Z/p\oplus\mathbb Z/p$: the latter is annihilated by $p$ and the former is not. There are no independent-event squares in a chain, so excision holds without imposing a split extension. This diagram is also the residual diagram of $\mathbb Z\xrightarrow{p}\mathbb Z\xrightarrow{p}\mathbb Z$.
\end{proof}

The replacement for ordered-filtration independence is therefore now explicit: the event poset controls which attachments depend on earlier events, excision controls independence from unrelated events, and the principal-ideal diagram retains the extension data needed for reconstruction.


\section{Obstruction theory}

Residual and obstruction should not be identified. A residual is a structured discrepancy. An obstruction is a discrepancy that cannot be removed by the admissible calculus.

\begin{definition}[Killable residual]
Let $\delta\in\Zoo(\Seeds_0)$. Call $\delta$ \emph{killable} if there exists an admissible finite constructor word $W$ such that
\[
W(\delta)\simeq0.
\]
\end{definition}

\begin{definition}[Persistent obstruction]
A residual seed $\delta$ is a \emph{persistent obstruction} relative to a constructor zoo if it is not killable by any admissible finite word.
\end{definition}

This definition is intentionally operational. More refined applications should replace it by an obstruction class in a quotient, derived functor, cohomology group, or localization.

A first taxonomy is:
\[
\boxed{
\Omega=(\Omega_R,\Omega_T,\Omega_{TR},\Omega_G)
}
\]
with:
\begin{itemize}
    \item $\Omega_R$: no admissible refinement exists;
    \item $\Omega_T$: a required within-resolution transformation does not exist;
    \item $\Omega_{TR}$: exchange cannot be made coherent or its residual persists;
    \item $\Omega_G$: locally compatible data fail to reconstruct globally.
\end{itemize}

\subsection{Universal obstruction quotient in the stable calculus}

Let $\mathcal Z$ be a small idempotent-complete stable $\infty$-category representing a thick seed zoo, and let $\mathscr M$ be a collection of exact constructor endofunctors of $\mathcal Z$ containing the identity and closed under composition.

\begin{definition}[Operationally killable subcategory]
Define
\[
\mathcal K_{\mathscr M}
:=
\{X\in\mathcal Z:\exists W\in\mathscr M,\;W(X)\simeq0\}.
\]
Let
\[
\mathcal N_{\mathscr M}
:=
\operatorname{Thick}(\mathcal K_{\mathscr M})
\]
be the smallest thick stable subcategory containing all operationally killable objects.
\end{definition}

The thick closure is necessary in general: two objects may be killed by unrelated constructor words with no single constructor word killing an extension of them.

\begin{definition}[Stable obstruction category and obstruction class]\label{def:obstruction-quotient}
The \emph{stable obstruction category} is the Verdier quotient
\[
\mathcal O_{\mathscr M}
:=
\mathcal Z/\mathcal N_{\mathscr M},
\]
with quotient functor
\[
q:\mathcal Z\longrightarrow\mathcal O_{\mathscr M}.
\]
For $X\in\mathcal Z$, define its \emph{obstruction class}
\[
\Obs_{\mathscr M}(X):=q(X).
\]
\end{definition}

\begin{theorem}[Universal stable obstruction quotient]\label{thm:obstruction-quotient}
The quotient $q$ is exact and satisfies:
\begin{enumerate}
    \item every operationally killable object maps to zero;
    \item $q(X)\simeq0$ if and only if $X\in\mathcal N_{\mathscr M}$;
    \item for every exact functor $E:\mathcal Z\to\mathcal D$ into a stable $\infty$-category such that $E(K)\simeq0$ for every $K\in\mathcal K_{\mathscr M}$, there is an exact factorization, unique up to contractible choice together with its comparison to $E$,
    \[
    \mathcal Z\xrightarrow{q}\mathcal O_{\mathscr M}\xrightarrow{\bar E}\mathcal D.
    \]
\end{enumerate}
Thus $\mathcal O_{\mathscr M}$ is universal among exact targets in which all operationally killable residuals vanish.
\end{theorem}

\begin{proof}
The Verdier quotient by the thick subcategory $\mathcal N_{\mathscr M}$ exists and the quotient functor is exact. Its kernel is exactly $\mathcal N_{\mathscr M}$. This gives (1)--(2). For (3), the kernel of an exact functor is thick. Since $E$ kills $\mathcal K_{\mathscr M}$, it kills the thick subcategory generated by that class, namely $\mathcal N_{\mathscr M}$. The universal property of the stable Verdier quotient gives the factorization. More precisely, precomposition with $q$ identifies the category of exact functors out of $\mathcal O_{\mathscr M}$ with the full subcategory of exact functors out of $\mathcal Z$ annihilating $\mathcal N_{\mathscr M}$. This also specifies the sense of uniqueness. See Drew~\cite{Drew2015} for the stable localization and \cite{StacksVerdier} for the kernel calculation in the triangulated homotopy category.
\end{proof}

\status{proved in the stable regime using the standard Verdier quotient}

A zero obstruction here means membership in the chosen thick kernel. It need not mean that one constructor word kills the object. Nor is $q(X)$ a chosen irreducible representative back inside $\mathcal Z$. A section or reflective localization would be extra data. The quotient also need not inherit the constructors automatically, as \cref{ex:quotient-collapse} will show.

\begin{example}[Why thick closure is genuinely necessary]\label{ex:killable-not-thick}
Let $\mathcal C$ be any nonzero stable category and set
\[
\mathcal Z=\mathcal C^3.
\]
Define exact idempotent endofunctors
\[
U(X_1,X_2,X_3)=(0,0,X_2\oplus X_3),
\qquad
V(X_1,X_2,X_3)=(0,0,X_1\oplus X_3).
\]
They satisfy
\[
U^2\simeq U,
\qquad
V^2\simeq V,
\qquad
U\circ V\simeq V,
\qquad
V\circ U\simeq U,
\]
so $\mathscr M=\{\id,U,V\}$ is closed under composition. Choose $0\not\simeq A\in\mathcal C$ and set
\[
X=(A,0,0),
\qquad
Y=(0,A,0).
\]
Then $U(X)\simeq0$ and $V(Y)\simeq0$, so both $X$ and $Y$ are operationally killable. However
\[
X\oplus Y=(A,A,0),
\]
and both $U$ and $V$ send this object to $(0,0,A)$, which is fixed by $U$ and $V$. Hence no constructor word kills $X\oplus Y$. Thus $\mathcal K_{\mathscr M}$ need not be additive, much less thick.
\end{example}

This example shows that replacing $\mathcal K_{\mathscr M}$ by its thick closure in \cref{def:obstruction-quotient} is forced if the quotient is required to be universal for exact functors.

\begin{definition}[Directed annihilator property]
The constructor system $\mathscr M$ has the \emph{directed annihilator property} if for every $U,V\in\mathscr M$ there exist $P,Q,W\in\mathscr M$ and natural equivalences
\[
W\simeq P\circ U\simeq Q\circ V.
\]
Thus any two constructor words admit a common postcomposition. With multiplication written as composition, $P\circ U=Q\circ V$ is a common left multiple. The displayed equation, rather than a left/right name, fixes the convention.
\end{definition}

\begin{theorem}[When operational killability is already thick]\label{thm:killable-thick}
If $\mathscr M$ consists of exact endofunctors and has the directed annihilator property, then $\mathcal K_{\mathscr M}$ is a thick stable subcategory. Hence
\[
\mathcal N_{\mathscr M}=\mathcal K_{\mathscr M}
\]
and
\[
\boxed{
\Obs_{\mathscr M}(X)\simeq0
\iff
X\text{ is operationally killable}.}
\]
\end{theorem}

\begin{proof}
The zero object is killed by every exact constructor. If $U(X)\simeq0$, exactness gives $U(X[n])\simeq U(X)[n]\simeq0$ for every integer $n$, so the class is closed under both suspensions and desuspensions. It is closed under retracts because exact functors preserve retract diagrams.

For cofibers, suppose $X$ is killed by $U$ and $Y$ is killed by $V$. Choose a common postcomposition
\[
W\simeq P\circ U\simeq Q\circ V.
\]
Then $W(X)\simeq0\simeq W(Y)$. For any morphism $f:X\to Y$, exactness gives
\[
W(\Cofib(f))\simeq\Cofib(Wf)\simeq\Cofib(0\to0)\simeq0.
\]
Thus the killable objects form a stable subcategory closed under retracts, hence a thick subcategory.
\end{proof}

\begin{example}[Idempotent projection model]
Let $\mathcal Z=\mathcal A\times\mathcal B$ for stable categories $\mathcal A,\mathcal B$, and let
\[
L(a,b)=(a,0).
\]
With constructor monoid generated by the idempotent exact functor $L$, an object is operationally killable exactly when it lies in $0\times\mathcal B$. The obstruction quotient is equivalent to $\mathcal A$, and
\[
\Obs(a,b)\simeq a.
\]
Thus the quotient literally retains the component that cannot be annihilated by the admissible constructor.
\end{example}

\begin{theorem}[Functorial transport of obstruction]\label{thm:obstruction-transport}
Let $(\mathcal Z_A,\mathscr M_A)$ and $(\mathcal Z_B,\mathscr M_B)$ be stable obstruction systems, and let
\[
\Lambda:\mathcal Z_A\longrightarrow\mathcal Z_B
\]
be exact. If
\[
\Lambda(\mathcal K_{\mathscr M_A})\subseteq\mathcal N_{\mathscr M_B},
\]
then $\Lambda$ induces an exact functor
\[
\bar\Lambda:\mathcal O_{\mathscr M_A}\longrightarrow\mathcal O_{\mathscr M_B}.
\]
If $\Lambda$ is an equivalence and carries $\mathcal N_{\mathscr M_A}$ onto $\mathcal N_{\mathscr M_B}$, then $\bar\Lambda$ is an equivalence. In that case obstruction vanishing is preserved and reflected.
\end{theorem}

\begin{proof}
The hypothesis and exactness imply
\[
\Lambda(\mathcal N_{\mathscr M_A})\subseteq\mathcal N_{\mathscr M_B},
\]
because exact functors preserve the operations used to generate a thick subcategory. Therefore $q_B\Lambda$ kills $\mathcal N_{\mathscr M_A}$, so the universal property in \cref{thm:obstruction-quotient} gives $\bar\Lambda$. If $\Lambda$ is an equivalence identifying the two thick kernels, a quasi-inverse descends in the same way and yields a quasi-inverse to $\bar\Lambda$.
\end{proof}

This theorem gives a precise version of the intended local-global principle: a correspondence need not merely preserve local invariants; it can transport the quotient in which removable residual structure has been annihilated.

\begin{theorem}[Obstruction localization along normalization]\label{thm:obstruction-localization}
Let
\[
X_0\xrightarrow{f_1}X_1\xrightarrow{}\cdots\xrightarrow{f_m}X_m
\]
be a normalization path as in \cref{thm:stable-filtration}, with elementary transported residuals $D_i=\Cofib(f_i)$ and total residual
\[
D_{\mathrm{tot}}=\Cofib(X_0\to X_m).
\]
Let
\[
q:\mathcal Z\to\mathcal O_{\mathscr M}
\]
be the stable obstruction quotient. Then:
\begin{enumerate}
    \item $q(D_{\mathrm{tot}})$ inherits a finite filtration whose successive cofibers are $q(D_i)$;
    \item if $q(D_i)\simeq0$ for every $i$, then
    \[
    q(D_{\mathrm{tot}})\simeq0;
    \]
    \item equivalently,
    \[
    q(D_{\mathrm{tot}})\not\simeq0
    \Longrightarrow
    \exists i,\;q(D_i)\not\simeq0.
    \]
\end{enumerate}
If the directed annihilator property of \cref{thm:killable-thick} holds, the phrase ``$q(D)\simeq0$'' can be replaced throughout by ``$D$ is operationally killable.''
\end{theorem}

\begin{proof}
The residual filtration in \cref{thm:stable-filtration} consists of cofiber sequences
\[
F_{i-1}\longrightarrow F_i\longrightarrow D_i,
\qquad
F_0\simeq0,
\qquad
F_m\simeq D_{\mathrm{tot}}.
\]
Because the Verdier quotient functor $q$ is exact, applying $q$ gives cofiber sequences
\[
q(F_{i-1})\longrightarrow q(F_i)\longrightarrow q(D_i).
\]
This proves (1). If every $q(D_i)$ vanishes, then $q(F_{i-1})\to q(F_i)$ is an equivalence for every $i$. Starting from $q(F_0)\simeq0$ yields $q(F_m)\simeq0$, proving (2). Statement (3) is its contrapositive. Under the directed annihilator hypothesis, \cref{thm:killable-thick} identifies quotient-vanishing with operational killability.
\end{proof}

\status{proved from the stable residual filtration and exactness of the obstruction quotient}

This result gives the calculus a precise diagnostic property: a persistent obstruction in a composite transformation-resolution word cannot appear from nowhere. Relative to a chosen normalization path, it must be visible in at least one elementary exchange residual after passage to the obstruction quotient.

\subsection{When the obstruction quotient must retain the constructors}

The universal property of $\mathcal Z/\mathcal N_{\mathscr M}$ concerns exact functors that kill a prescribed class. It does not automatically preserve the full transformation grammar. To make the quotient a new environment for that grammar, the kernel must be invariant under the constructors.

\begin{example}[A constructor can fail to descend, and invariant closure can collapse]\label{ex:quotient-collapse}
Return to \cref{ex:killable-not-thick}, with $\mathcal Z=\mathcal C^3$ and $\mathscr M=\{\id,U,V\}$. An object is killed by $U$ precisely when its second and third components vanish, and killed by $V$ precisely when its first and third components vanish. Indeed, a summand of a zero direct sum is zero. Consequently
\[
\mathcal N_{\mathscr M}=\mathcal C\times\mathcal C\times0,
\qquad
\mathcal Z/\mathcal N_{\mathscr M}\simeq\mathcal C.
\]
Nevertheless
\[
U(0,A,0)=(0,0,A)
\]
for every $A\in\mathcal C$. The input lies in $\mathcal N_{\mathscr M}$ and the output generally does not. If $U$ descended to an exact $\bar U$ on the quotient, then
$qU(0,A,0)\simeq\bar Uq(0,A,0)\simeq0$, a contradiction for $A\not\simeq0$.

Moreover, any constructor-invariant thick subcategory containing all killable objects contains both the first and second coordinate subcategories. Applying $U$ to the second coordinate forces it to contain the third. Closure under direct sums then gives all of $\mathcal Z$. Thus the constructor-compatible quotient in this example is the zero category.
\end{example}

\begin{definition}[Constructor-invariant thick saturation]\label{def:invariant-saturation}
For a small stable constructor system $(\mathcal Z,\mathscr M)$ as above, put
\[
\mathcal J_{\mathscr M}
:=\operatorname{Thick}\{W(K):W\in\mathscr M,
\ K\in\mathcal K_{\mathscr M}\}.
\]
Its quotient is denoted
\[
q^{\mathrm{inv}}:\mathcal Z\longrightarrow
\mathcal O^{\mathrm{inv}}_{\mathscr M}
:=\mathcal Z/\mathcal J_{\mathscr M}.
\]
The identity constructor ensures $\mathcal N_{\mathscr M}\subseteq\mathcal J_{\mathscr M}$.
\end{definition}

\begin{theorem}[Universal constructor-compatible obstruction quotient]\label{thm:invariant-quotient}
The subcategory $\mathcal J_{\mathscr M}$ is the least thick subcategory that contains $\mathcal K_{\mathscr M}$ and is preserved by every constructor. Each $T\in\mathscr M$ descends to an exact endofunctor of $\mathcal O^{\mathrm{inv}}_{\mathscr M}$. If the constructor action is supplied with coherent unit and composition data, those data descend as well.

The quotient is universal among exact functors that annihilate every object $W(K)$ in the displayed generating family. In particular, every exact realization that annihilates killable objects and intertwines the given constructors with exact constructors of its target factors through $q^{\mathrm{inv}}$.
\end{theorem}

\begin{proof}
Fix $T\in\mathscr M$. It takes a generator $W(K)$ to $(T\circ W)(K)$, again a generator because $\mathscr M$ is closed under composition. Exactness and preservation of retracts imply
\[
T(\mathcal J_{\mathscr M})\subseteq\mathcal J_{\mathscr M}.
\]
Conversely any invariant thick subcategory containing $\mathcal K_{\mathscr M}$ contains every $W(K)$ and therefore their thick closure. This proves minimality.

Now $q^{\mathrm{inv}}T$ annihilates $\mathcal J_{\mathscr M}$, so the quotient universal property gives $\bar T$ with
\[
\bar T q^{\mathrm{inv}}\simeq q^{\mathrm{inv}}T.
\]
The full faithfulness of precomposition in the functor-category universal property descends any supplied unit, composition, and higher coherence data. It does not create action data absent from the original system.

The last universal property follows directly from thick closure and Verdier localization. If $E$ is equivariant and kills $K$, then
$E(WK)\simeq W'E(K)\simeq W'(0)\simeq0$. Thus it kills all the generators and factors through the quotient.
\end{proof}

\status{proved for small stable exact constructor systems}

There is no general nontriviality conclusion. The zero category is a legitimate universal quotient. A useful application must show that the chosen admissible repairs do not annihilate everything after invariant thick closure.

\begin{theorem}[Directed annihilators also give constructor invariance]\label{thm:directed-invariant}
Under the directed annihilator hypothesis of \cref{thm:killable-thick},
\[
\mathcal J_{\mathscr M}=\mathcal N_{\mathscr M}
=\mathcal K_{\mathscr M}.
\]
In particular the same quotient both supports the admitted constructors and detects operational killability exactly.
\end{theorem}

\begin{proof}
The earlier theorem shows that $\mathcal K_{\mathscr M}$ is thick. It remains to prove invariance. Suppose $U(X)\simeq0$ and fix a constructor $T$. Directedness supplies constructors $P,Q,W$ with
\[
W\simeq P\circ U\simeq Q\circ T.
\]
Therefore
\[
Q(TX)\simeq P(UX)\simeq P(0)\simeq0.
\]
Thus $TX$ is killable. The killable subcategory is invariant and thick, so minimality of $\mathcal J_{\mathscr M}$ gives the asserted equalities.
\end{proof}

\begin{proposition}[A noncollapse certificate]\label{prop:noncollapse}
Suppose $E:\mathcal Z\to\mathcal D$ is an exact equivariant realization into a stable constructor system, $E$ kills every operationally killable object, and $E(X)\not\simeq0$ for some $X$. Then
\[
q^{\mathrm{inv}}(X)\not\simeq0,
\qquad
\mathcal O^{\mathrm{inv}}_{\mathscr M}\not\simeq0.
\]
\end{proposition}

\begin{proof}
By \cref{thm:invariant-quotient}, write $E\simeq\bar E q^{\mathrm{inv}}$. Exactness of $\bar E$ would imply $E(X)\simeq0$ if $q^{\mathrm{inv}}(X)\simeq0$, a contradiction.
\end{proof}

For the projection example $\mathcal Z=\mathcal A\times\mathcal B$ and $L(a,b)=(a,0)$, projection to $\mathcal A$ is such a detector whenever $\mathcal A$ is nonzero. This gives a nontrivial model, in contrast with \cref{ex:quotient-collapse}. Admissibility matters: allowing the zero functor as a repair makes every object killable and necessarily collapses the quotient.

The refined distinction is therefore
\[
\begin{aligned}
\mathcal K_{\mathscr M}&:\ \text{objects killed by a single admitted word},\\
\mathcal N_{\mathscr M}&:\ \text{their thick closure},\\
\mathcal J_{\mathscr M}&:\ \text{their constructor-invariant thick closure}.
\end{aligned}
\]
Only under additional hypotheses can these three classes be identified.


\begin{problem}[Nonstable intrinsic obstruction]
Construct an analogue of \cref{def:obstruction-quotient} before stable/additive completion, or prove that the stable quotient is the correct universal recipient for the applications of interest.
\end{problem}

\section{Universal excisive localization}\label{sec:excisive-localization}

The constructor-invariant quotient can now be attached to a precise target property: making the event diagram excisive. This is different from declaring every residual removable.

Let $\mathcal Z$ be a small stable category, let $\mathscr M$ be a coherently acting monoid of exact endofunctors, and let $\mathscr X$ be a small family of finite coherent event diagrams in $\mathcal Z$. Denote by $\mathscr H$ the family of interaction residuals of all their elementary independent-event squares. Put
\[
 \mathcal J_{\square}
 =\operatorname{Thick}\{T(H):T\in\mathscr M,\ H\in\mathscr H\}.
\]
This notation does not assert that $H$ was operationally killable before localization.

\begin{theorem}[Universal constructor-compatible excisive quotient]\label{thm:excisive-quotient}
The stable Verdier quotient
\[
 q_{\square}:\mathcal Z\longrightarrow\mathcal Q_{\square}
 :=\mathcal Z/\mathcal J_{\square}
\]
has the following properties.
\begin{enumerate}
\item The constructors descend, with their supplied coherent action data.
\item Every diagram $q_{\square}X$, for $X\in\mathscr X$, has context-independent event residuals and satisfies principal-ideal reconstruction.
\item Every exact equivariant functor $E:\mathcal Z\to\mathcal D$ for which all the diagrams $EX$ are excisive factors through $q_{\square}$, with the enhanced universal uniqueness of Verdier localization.
\item The kernel is the least constructor-invariant thick subcategory forcing those squares to be excisive.
\end{enumerate}
The quotient is not automatically nonzero. If such an $E$ has $E(Y)\not\simeq0$, then $q_{\square}(Y)\not\simeq0$.
\end{theorem}

\begin{proof}
The proof of \cref{thm:invariant-quotient} applies to any designated family of objects, not only killable objects. A constructor sends a generator $T(H)$ to another generator. Exactness then preserves its thick closure. The Verdier universal property therefore descends constructors and their already specified coherence.

Since $q_{\square}(H)\simeq0$ and exact functors preserve interaction residuals, every prescribed elementary square in $q_{\square}X$ is a pushout. Apply \cref{thm:event-context-invariance,thm:principal-reconstruction} to obtain (2).

If $EX$ is excisive, then $E(H)\simeq0$. Equivariance implies $E(T(H))\simeq T'(E(H))\simeq0$. Exactness annihilates the thick closure, so $E$ factors through the quotient. Conversely, the kernel of any exact equivariant localization making the diagrams excisive must contain $\mathscr H$ and its invariant thick closure. This proves the universal and minimality assertions. Finally a factorization $E\simeq\bar E q_{\square}$ would send $Y$ to zero if $q_{\square}(Y)$ were zero.
\end{proof}

The external input is stable Verdier localization, as in \cite{Drew2015,StacksVerdier}. The new use here is the explicitly chosen family of square interaction residuals. It does not supply an arithmetic interpretation for a quotient unless that interpretation is proved separately.

\section{An integral seed model and prime-overlap calculus}\label{sec:integral-model}

We now give a completely specified arithmetic instance. It uses perfect complexes over $\mathbb Z$, not Galois representations or automorphic representations. Its purpose is to make the coherence, residual, excision, noncollapse, and finite-generator obligations explicit in one model.

\subsection{Tensor cubes and iterated residuals}

\begin{theorem}[Tensor-cube residual formula]\label{thm:tensor-cube}
Let $\mathcal S$ be a symmetric monoidal stable category whose tensor product is exact separately in each variable. For $m\ge1$ and arrows $f_i:A_i\to B_i$, form the coherent $m$-cube
\[
 Y(S)=\bigotimes_{i=1}^m Z_i(S),\qquad
 Z_i(S)=\begin{cases}B_i&i\in S,\\A_i&i\notin S.\end{cases}
\]
Its total cofiber, obtained by taking the cofiber in every coordinate, satisfies
\[
 \operatorname{tcofib}(Y)\simeq
 \bigotimes_{i=1}^m\Cofib(f_i).
\]
Equivalently it is the cofiber of the canonical map from the colimit of the punctured cube to its terminal vertex. The equivalences are natural in the arrows, with the coherent symmetries of the monoidal structure, rather than literal equality of chain matrices for different coordinate orders.
\end{theorem}

\begin{proof}
For $m=1$ this is the definition. Taking cofibers in the last coordinate replaces its factor $A_m\to B_m$ by $\Cofib(f_m)$ because tensoring by the other factors is exact. Repeat for the remaining coordinates. Interchanging finite colimits shows that this iterated construction is the cofiber of the punctured-cube colimit map. The universal properties give naturality and coherent interchange. In a complex model, these interchange maps incorporate the tensor differential's graded signs.
\end{proof}

For a commutative ring $R$ and $a\in R$, set
\[
 C_R(a)=\bigl[R\xrightarrow{a}R\bigr],
\]
with the source in cohomological degree $-1$ and the target in degree $0$. The tensor product $\bigotimes_i C_R(a_i)$ is the Koszul complex in this cohomological convention. Its tensor, cone, and change-of-basis properties are classical; see \cite{StacksKoszul}. We use them with an explicit sign convention and direct proofs below.

\subsection{Scalar exchange diagrams}

Fix a finite poset $P$ and assign a positive integer $n_e$ to every event $e\in P$. For $M\in\operatorname{Perf}(\mathbb Z)$ define
\[
 X_M(I)=M,\qquad
 X_M(I\subseteq J)=\left(\prod_{e\in J\setminus I}n_e\right)\id_M.
\]
Empty products are $1$. For nested ideals these products multiply exactly, so $X_M$ is a strict chain-level diagram followed by the stable realization. When $P=P_W$, take identity transformation and coarsening functors and let the exchange cell for event $e$ be multiplication by $n_e$. This realizes the diagram inside the typed two-axis grammar of \cref{thm:strict-event-realization}. The object at every vertex is the same, but its comparison maps retain nontrivial arithmetic.

For an enabled event $e$,
\[
 E_e(I)\simeq M\otimes^{\mathbf L}_{\mathbb Z}C_{\mathbb Z}(n_e).
\]
The comparison after an unrelated event $f$ is multiplication by $n_f$ on that complex. Consequently
\[
 \boxed{\Theta_{e,f}(I)\simeq
 M\otimes^{\mathbf L}_{\mathbb Z}
 C_{\mathbb Z}(n_e)\otimes^{\mathbf L}_{\mathbb Z}C_{\mathbb Z}(n_f).}
\]
The complexes $C_{\mathbb Z}(n)$ are bounded and free, so ordinary tensor complexes compute these derived tensors. The total normalization residual is instead
\[
 \mathcal F_{X_M}(P)\simeq
 M\otimes^{\mathbf L}_{\mathbb Z}C_{\mathbb Z}(N),
 \qquad N=\prod_{e\in P}n_e.
\]
This is not the total cofiber of the entire event cube. The normalization residual compares the initial and final vertices, whereas a cube's total cofiber measures their higher interaction. These are different constructions.

\begin{theorem}[Integral interaction is supported at the gcd]\label{thm:integral-gcd}
For positive integers $a_1,\ldots,a_m$, $m\ge1$, put $g=\gcd(a_1,\ldots,a_m)$. Then
\[
 K(a_1,\ldots,a_m):=
 \bigotimes_{i=1}^m C_{\mathbb Z}(a_i)
 \simeq\bigoplus_{j=0}^{m-1}
 C_{\mathbb Z}(g)[j]^{\oplus\binom{m-1}{j}}.
\]
The displayed splitting depends on a unimodular change of basis and is not claimed canonical. Therefore
\[
 H^{-j}K(a_1,\ldots,a_m)\cong
 (\mathbb Z/g)^{\oplus\binom{m-1}{j}}
 \quad(0\le j\le m-1),
\]
and every other cohomology group vanishes. In particular the interaction vanishes exactly when $g=1$.
\end{theorem}

\begin{proof}
The Euclidean algorithm produces a unimodular matrix $U\in\operatorname{GL}_m(\mathbb Z)$ with
\[
 (a_1,\ldots,a_m)U=(g,0,\ldots,0).
\]
The Koszul complex is the exterior algebra on a free module in degree $-1$, with differential the derivation determined by this row vector. Hence the maps $\bigwedge^jU$ give an isomorphism of complexes from $K(g,0,\ldots,0)$ to $K(a_1,\ldots,a_m)$.

Each zero entry contributes $C_{\mathbb Z}(0)\cong\mathbb Z\oplus\mathbb Z[1]$ with zero differential. Distributing these $m-1$ factors gives the stated shifted copies, with the standard shift isomorphisms. Since $g>0$, the map $g:\mathbb Z\to\mathbb Z$ is injective. Thus $C_{\mathbb Z}(g)$ has only $H^0=\mathbb Z/g$, proving the cohomology formula and the vanishing criterion.
\end{proof}

For two events the entire calculation is visible:
\[
 K(a,b)=\left[
 \mathbb Z\xrightarrow{\binom{-b}{a}}
 \mathbb Z^2\xrightarrow{(a\ \ b)}\mathbb Z
 \right],
 \qquad\text{degrees }-2,-1,0.
\]
The product of the two differentials is zero. The kernel of $(a\ b)$ is generated by $(-b/g,a/g)$, and the incoming image is $g$ times that generator. Both $H^{-1}$ and $H^0$ are $\mathbb Z/g$. Equivalently, if $au+bv=g$, then
\[
 U=\begin{pmatrix}u&-b/g\\v&a/g\end{pmatrix},
 \qquad\det U=1,
 \qquad(a\ b)U=(g\ 0).
\]

Thus a strictly commuting scalar exchange square is a homotopy pushout on the seed $\mathbb Z$ exactly when the two exchange integers are coprime. Its commutativity alone says nothing about this condition.

\begin{example}[Higher vanishing does not remove lower interactions]\label{ex:higher-not-lower}
For $(6,10,15)$, all pair interactions are nonzero: their gcds are $2,3,5$. But the threefold interaction is zero because the joint gcd is $1$. Hence vanishing of a higher total cofiber does not imply vanishing of the lower square interactions. It cannot replace the square-by-square excision condition in \cref{thm:event-context-invariance}.
\end{example}

\subsection{A nontrivial universal quotient with an exact detector}

Let $\mathscr M$ consist of tensoring by perfect complexes, closed under composition with its coherent tensor action. This is a constructor system, not the repair monoid defining operational killability in the earlier sections. We now kill only the designated interaction residuals.

For the fixed labelled poset define the finite prime set
\[
 S_P=\bigcup_{e\parallel f}
 \{p\text{ prime}:p\mid\gcd(n_e,n_f)\},
\]
where $e\parallel f$ means incomparable. Let $\mathcal J_P$ be the invariant thick closure of the square interaction residuals for the unit seed $\mathbb Z$.

\begin{theorem}[Prime-overlap classification of the excisive kernel]\label{thm:prime-overlap-quotient}
In $\operatorname{Perf}(\mathbb Z)$,
\[
 \boxed{\mathcal J_P=
 \operatorname{Thick}\{\mathbb Z/p:p\in S_P\}.}
\]
Equivalently, $\mathcal J_P$ consists exactly of perfect complexes whose cohomology groups are finite abelian groups with prime divisors in $S_P$. For the stable Verdier quotient $q_P$, and every perfect complex $M$,
\[
 q_P(M)\simeq0
 \quad\Longleftrightarrow\quad
 M\otimes^{\mathbf L}_{\mathbb Z}\mathbb Z[S_P^{-1}]\simeq0.
\]
The quotient is nonzero. Its constructor action is defined, every $q_PX_M$ is excisive, and its reduced residual diagram is reconstructed from principal ideals.
\end{theorem}

\begin{proof}
First, every pair of incomparable events occurs in an elementary square. Take the ideal $D_e\cup D_f$. It contains neither $e$ nor $f$, contains their predecessors, and enables both. The square residual at the unit is $K(n_e,n_f)$.

By \cref{thm:integral-gcd}, this complex has the same thick closure as $C_{\mathbb Z}(g)$ for $g=\gcd(n_e,n_f)$: it is a direct sum of that complex and its shift, and contains it as a retract. The thick closure of any family in $\operatorname{Perf}(\mathbb Z)$ is already a tensor ideal. Indeed $\operatorname{Perf}(\mathbb Z)=\operatorname{Thick}(\mathbb Z)$, so tensoring a fixed generator by a perfect complex builds its image using finite cones, shifts, and retracts from that same generator.

For $p\mid g$, tensor $C_{\mathbb Z}(g)$ by $C_{\mathbb Z}(p)$. The gcd theorem identifies the result with $C_{\mathbb Z}(p)\oplus C_{\mathbb Z}(p)[1]$. Hence $\mathbb Z/p\simeq C_{\mathbb Z}(p)$ belongs to the thick closure. Conversely, the Chinese remainder decomposition of $\mathbb Z/g$ and the short exact sequences building $\mathbb Z/p^j$ from copies of $\mathbb Z/p$ put $C_{\mathbb Z}(g)\simeq\mathbb Z/g$ in the thick subcategory generated by its prime divisors. This proves the boxed equality.

Having finite $S_P$-primary cohomology is preserved by finite triangles and retracts, by the long exact sequence in cohomology. Conversely, each such finite abelian cohomology group is built by finite extensions from the indicated $\mathbb Z/p$. A perfect complex is bounded, and its finite truncation triangles then build it from these groups. These truncations are perfect over $\mathbb Z$: every finitely generated abelian group has a finite free resolution of length at most one. This proves the cohomological characterization.

Localization of abelian groups is exact. Thus tensoring with $\mathbb Z[S_P^{-1}]$ has cohomology obtained by localizing each cohomology group. A finitely generated abelian group localizes to zero exactly when it is finite with all prime divisors in $S_P$. The localization kernel is therefore $\mathcal J_P$, as is the kernel of the Verdier quotient by this thick subcategory. This proves the equivalence of vanishing tests on every source object.

Finally $\mathbb Z[S_P^{-1}]$ is a nonzero ring, so localization detects that $q_P(\mathbb Z)\not\simeq0$. It is exact and intertwines tensor constructors with their localized counterparts. The universal excisive-quotient theorem supplies the remaining assertions. For general $M$, each square residual is its unit-seed counterpart tensored with $M$, so it lies in the same invariant kernel.
\end{proof}

The positivity hypothesis excludes a real collapse case. If two independent labels are both zero, then $K(0,0)\simeq\mathbb Z\oplus\mathbb Z[1]^{\oplus2}\oplus\mathbb Z[2]$ contains the unit as a retract. Killing that interaction forces the whole perfect category to vanish. Thus noncollapse is proved for the specified positive-label model, not for arbitrary scalar labels.

No equivalence of the entire quotient category with a category of localized perfect complexes is needed for this theorem. The proved statement is equality of kernels, a nonzero exact detector, and the universal excisive property. This avoids inserting an unproved full-faithfulness or idempotent-completion claim.

\begin{example}[Remove shared torsion without deleting the primary residuals]\label{ex:six-ten}
Take two independent events labelled $6$ and $10$. Their interaction is
\[
 K(6,10)\simeq\mathbb Z/2\oplus(\mathbb Z/2)[1].
\]
Here $S_P=\{2\}$. Inverting $2$ kills this interaction while retaining
\[
 C_{\mathbb Z}(6)\otimes\mathbb Z[1/2]\simeq\mathbb Z/3,
 \qquad
 C_{\mathbb Z}(10)\otimes\mathbb Z[1/2]\simeq\mathbb Z/5.
\]
The total normalization residual is $C_{\mathbb Z}(60)$. After localization it is $\mathbb Z/15$, and principal-ideal reconstruction gives
\[
 \mathbb Z/15\simeq\mathbb Z/3\oplus\mathbb Z/5.
\]
This is a classical Chinese-remainder instance of the calculus. Before localization the two paths have first residuals $\mathbb Z/6$ and $\mathbb Z/10$. After localization those first stages remain different, so the original ordered-filtration independence claim is still not being restored.
\end{example}

\subsection{A finite certificate for an infinite scalar grammar}

\begin{corollary}[Finite-signature arithmetic certification]\label{cor:finite-signature}
Fix finitely many positive integer exchange labels $n_1,\ldots,n_t$ and a specified symmetric relation $\mathcal I$ of label pairs allowed to occur on incomparable events. The relation may contain a diagonal pair. Put
\[
 S_{\mathcal I}=\bigcup_{(i,j)\in\mathcal I}
 \{p:p\mid\gcd(n_i,n_j)\}.
\]
Killing the finitely many unit-seed interactions $K(n_i,n_j)$ for $(i,j)\in\mathcal I$, with invariant thick closure, makes every scalar event diagram on every finite labelled poset satisfying this independence restriction excisive, for every perfect input $M$. The quotient is nonzero and its kernel is the finite-prime torsion class supported on $S_{\mathcal I}$.
\end{corollary}

\begin{proof}
There are only finitely many prescribed pair types. Every elementary square in every admitted diagram has one of these types, and its residual is the corresponding unit-seed Koszul complex tensored with $M$. All such residuals lie in the same invariant thick kernel. Apply \cref{thm:excisive-quotient,thm:prime-overlap-quotient}. The latter proof depends only on the finite family of gcds, so it applies unchanged to the signature-wide family.
\end{proof}

This is a genuine finite reduction of an unbounded word-length obligation, but only in the scalar integral regime with the stated independence relation. If a label can occur independently of itself, its primes belong to $S_{\mathcal I}$ and its primary torsion is erased. Declaring independence therefore changes the quotient's information content. Dependent events can share a prime without forcing its removal, as the $p$-then-$p$ chain in \cref{cor:independent-splitting} demonstrates.

The construction resolves a concrete semantic and obstruction problem in arithmetic. It does not identify these scalar kernels with automorphic occupancy, a ramification obstruction, or a global Langlands parameter.

\subsection{Beyond scalar exchanges: an integral matrix criterion}

The noncollapse test extends to tensor cubes of finite integer matrices. This is a different concrete realization from identity functors with scalar exchange cells: its independent directions are tensor factors of actual two-term complexes.

\begin{theorem}[Finite integral matrix collapse criterion]\label{thm:matrix-collapse}
Let $f_i:\mathbb Z^{a_i}\to\mathbb Z^{b_i}$ be finitely many integer matrices. Specify which pairs of directions may occur independently. For each matrix fix a Smith normal form, with rank $r_i$ and positive nonzero diagonal entries $d_{i,1},\ldots,d_{i,r_i}$. Put
\[
 \epsilon_i=\begin{cases}
 1&a_i+b_i-2r_i>0,\\
 0&a_i+b_i-2r_i=0,
 \end{cases}
 \qquad
 P_i=\{p:p\mid d_{i,k}\text{ for some }k\}.
\]
Let $\mathcal J$ be the tensor-invariant thick subcategory generated by $\Cofib(f_i)\otimes^{\mathbf L}\Cofib(f_j)$ over the allowed independent pairs.
\begin{enumerate}
\item If an allowed pair has $\epsilon_i=\epsilon_j=1$, then $\mathcal J=\operatorname{Perf}(\mathbb Z)$ and its quotient is zero.
\item Otherwise the quotient is nonzero and
\[
 \mathcal J=\operatorname{Thick}\{\mathbb Z/p:p\in S\},
\]
where $S$ is the union over allowed pairs of
\[
 S_{ij}=\begin{cases}
 P_i\cap P_j&\epsilon_i=\epsilon_j=0,\\
 P_j&\epsilon_i=1,\ \epsilon_j=0,\\
 P_i&\epsilon_i=0,\ \epsilon_j=1.
 \end{cases}
\]
Localization at $S$ is an exact nonzero detector with that kernel. It makes the prescribed tensor-cube squares excisive.
\end{enumerate}
A supplied Smith form is a finite certificate: its row and column changes must be unimodular and their matrix product must give the claimed diagonal form. The existence of Smith forms over $\mathbb Z$ follows from the Euclidean algorithm; the argument below needs only such a verified form.
\end{theorem}

\begin{proof}
Changing bases in the source and target yields a chain isomorphism
\[
 \Cofib(f_i)\simeq
 \mathbb Z^{b_i-r_i}\oplus\mathbb Z^{a_i-r_i}[1]
 \oplus\bigoplus_{k=1}^{r_i}C_{\mathbb Z}(d_{i,k}).
\]
Unit diagonal entries contribute contractible complexes. Thus $\epsilon_i=1$ means that a shift of the unit occurs as a direct summand. If two independent cones each have such a summand, their tensor product contains a shift of the unit. A thick subcategory containing it is all of $\operatorname{Perf}(\mathbb Z)$.

If no allowed pair has two free summands, every generating tensor is torsion. A free summand tensored with a torsion cone preserves that torsion's prime support. Two torsion summands have tensor product $K(d_{i,k},d_{j,l})$, whose thick support is the prime intersection by \cref{thm:integral-gcd}. Taking finite direct sums gives exactly the listed sets $S_{ij}$. The argument of \cref{thm:prime-overlap-quotient} identifies the generated thick subcategory with finite $S$-primary cohomology and identifies its kernel with localization at $S$. The localized unit is nonzero. The tensor-cube formula and the excisive-quotient theorem finish the proof.
\end{proof}

This criterion shows why noncollapse cannot be inferred from commutativity. Two commuting tensor directions with free cofiber components can force total collapse when their interaction is declared zero. Pure torsion interactions, in contrast, admit the precise finite-prime localization just proved.


\section{Arithmetic resolution towers and nilpotent interactions}\label{sec:arithmetic-towers}

The integral model of \cref{sec:integral-model} has nontrivial exchange maps but identity coarsening functors. We now replace those coarsenings by derived reduction between different coefficient rings. Three questions must be kept separate: whether an exchange is canonically exact, what a specified noninvertible exchange cell leaves behind, and whether forcing that residual to vanish destroys the coefficient category.

\subsection{Derived reduction is a genuine resolution operation}

Fix a finite nonempty set $\mathcal P$ of primes and put
\[
 D=\prod_{p\in\mathcal P}p,\qquad B_n=\mathbb Z/D^n\mathbb Z\quad(n\geq1).
\]
Write $B_{p,n}=\mathbb Z/p^n\mathbb Z$. The Chinese remainder isomorphisms
\[
 B_n\cong\prod_{p\in\mathcal P}B_{p,n}
\]
are compatible with reduction. Use the enhanced category $\mathcal S_n=\operatorname{Perf}(B_n)$ and the coarsening
\[
 C_{m,n}(M)=M\otimes^{\mathbf L}_{B_m}B_n,\qquad m\geq n.
\]
Derived tensor associativity supplies coherent equivalences $C_{l,n}\simeq C_{m,n}C_{l,m}$. These are the data of an enhanced indexed system. Literal equalities are not asserted between arbitrarily chosen derived-functor models. On specified bounded projective complexes the operation is termwise reduction. Preservation of perfect complexes under derived base change is standard \cite{StacksPerfect}; the pointwise cone model also shows exactness directly.

\begin{theorem}[Nonidentity conservative resolution tower]\label{thm:adic-tower}
The functors $C_{m,n}$ are exact and conservative on perfect complexes. They need not be faithful. In particular, for $m>n$, the nonzero endomorphism $D^n:B_m\to B_m$ becomes zero under $C_{m,n}$.
\end{theorem}
\begin{proof}
Exactness follows by reducing bounded projective cone diagrams. For conservativity it suffices, by the Chinese remainder decomposition and composition with reduction to level one, to prove that a perfect complex over $R=\mathbb Z/p^m$ is zero if its derived reduction to $\mathbb F_p$ is zero.

A finite projective module over $R$ is free. Indeed, lift a basis modulo $(p)$ to obtain a surjection $R^r\to P$. It is surjective because its cokernel $Q$ satisfies $Q=pQ$, whence $Q=p^mQ=0$. The surjection splits, and its kernel has zero reduction modulo $p$, so the same argument kills the kernel. A perfect complex therefore has a bounded finite free representative.

If a differential has a unit entry, elementary changes of bases split off the contractible two-term complex $R\xrightarrow{1}R$. The adjacent differentials have zero components into and out of that summand because their composites are zero. Repeating strictly decreases the total rank and terminates in a complex all of whose differential entries lie in $(p)$. Its reduction has zero differentials. If that reduction is acyclic, every remaining rank is zero. The original complex is contractible. Applied to a cone, this also proves that reduction reflects equivalences.

Finally $D^n$ is not zero modulo $D^m$, but is zero modulo $D^n$. Since $\operatorname{Hom}_{D(B_m)}(B_m,B_m)=B_m$, this is genuine loss of morphism information, not just a change of representatives.
\end{proof}

Conservativity detects zero objects and equivalences. It does not recover every morphism, a preferred lift, or the full information in a fine presentation.

\begin{example}[Why ordinary tensor and cohomology truncation are not substitutes]\label{ex:derived-not-naive}
Over $R=\mathbb Z/p^2$, the exact sequence $0\to(p)\to R\to\mathbb F_p\to0$ ceases to be left exact after ordinary tensor with $\mathbb F_p$: the first map becomes zero. The module $\mathbb F_p$ has the infinite periodic free resolution with every differential multiplication by $p$. Tensoring that resolution with $\mathbb F_p$ gives zero differentials, hence $\operatorname{Tor}_i^R(\mathbb F_p,\mathbb F_p)=\mathbb F_p$ for every $i\geq0$.

In contrast, $C_R(p)=[R\xrightarrow p R]$ is perfect, but has both $H^{-1}=\mathbb F_p$ and $H^0=\mathbb F_p$. It is not the module $\mathbb F_p$ placed in degree zero. We retain the actual two-term complex throughout the tower. Derived tensor is computed on these projective representatives, not on a truncated cohomology surrogate \cite{StacksDerivedTensor}.
\end{example}

\subsection{Coherent cells and their actual source of residuals}

For a fixed perfect integral complex $L$, put
\[
 A_{L,n}(M)=(L\otimes^{\mathbf L}_{\mathbb Z}B_n)\otimes^{\mathbf L}_{B_n}M.
\]
Base change gives a natural equivalence
\[
 \theta_{L;m,n}:A_{L,n}C_{m,n}\simeq C_{m,n}A_{L,m}.
\]
Thus the canonical tensor/base-change exchange already has zero residual. The tower alone does not force a defect. To test a noninvertible exchange, we explicitly add a positive integer label $b$ and use the cell
\[
 \eta_{L,b;m,n}=b^{m-n}\theta_{L;m,n}.
\]
For scalar unit templates take $L=\mathbb Z$, so that the transformation functors are identities but the coarsening functors are not.

\begin{theorem}[Coherent weighted exchange and residual transport]\label{thm:weighted-adic-exchange}
The cells $\eta_{L,b;m,n}$ are natural and coherent with composition of resolution changes. Their residual at $M\in\mathcal S_m$ is
\[
 \Cofib(\eta_{L,b;m,n,M})\simeq
 A_{L,n}C_{m,n}(M)\otimes^{\mathbf L}_{B_n}C_{B_n}(b^{m-n}).
\]
For a fixed scalar $a$, subsequent base change transports its cone exactly:
\[
 C_{n,k}\bigl(N\otimes^{\mathbf L}_{B_n}C_{B_n}(a)\bigr)
 \simeq C_{n,k}(N)\otimes^{\mathbf L}_{B_k}C_{B_k}(a).
\]
At a common target level $n$, two independent scalar events of positive lengths $d,e$ and labels $a,b$ have interaction
\[
 N\otimes^{\mathbf L}_{B_n}K_{B_n}(a^d,b^e),\qquad
 K_R(u,v)=C_R(u)\otimes^{\mathbf L}_R C_R(v).
\]
These formulas also hold with compatible tensor templates absorbed into $N$.
\end{theorem}
\begin{proof}
Naturality follows from central scalar multiplication. For $l\geq m\geq n$, the two-step cell has scalar $b^{m-n}b^{l-m}=b^{l-n}$. The associativity coherence of derived tensor identifies its underlying $\theta$ with $\theta_{L;l,n}$. Identity intervals have scalar one. Tensoring transformation templates gives their usual associative action, with scalar weights multiplied; no additional cells are inferred without these specified data.

After using $\theta$ to identify the target, the exchange is multiplication by $b^{m-n}$ on $A_{L,n}C_{m,n}M$. Exact tensor identifies its cone with the first formula. Associativity of derived tensor proves the second. The two-variable interaction formula follows by taking the two cofibers, as in \cref{thm:tensor-cube}. For a pure adjacent-swap word, tensor templates keep their relative order and all coarsenings combine into the same final reduction. Identifying those endpoints by the supplied tensor associativities leaves precisely the products of the scalar cells. This constructs the corresponding coherent scalar event diagram at the target level.
\end{proof}

\begin{remark}[Specified cells, not a claimed arithmetic correspondence]
Nonzero residuals in this model come from the chosen noninvertible scalar cells. They are not defects of the canonical base-change equivalence, and they are not identified with a ramification or automorphy obstruction. The construction supplies a reproducible nonidentity resolution test bed with explicit semantics.
\end{remark}

\subsection{The finite-level Koszul calculation}

\begin{theorem}[Prime-power interaction cohomology]\label{thm:prime-power-koszul}
Let $R=\mathbb Z/p^n$, let $a_1,\ldots,a_t$ be positive integers with $t\geq1$, and put
\[
 h=\min\{n,v_p(a_1),\ldots,v_p(a_t)\},\qquad Q_h=R/(p^h).
\]
Here $Q_0=0$ and $Q_n=R$. There is a generally noncanonical chain isomorphism
\[
 K_R(a_1,\ldots,a_t)\cong
 \bigoplus_{j=0}^{t-1}C_R(p^h)[j]^{\oplus\binom{t-1}{j}}.
\]
Consequently
\[
 \boxed{H^{-j}K_R(a_1,\ldots,a_t)\cong Q_h^{\oplus\binom tj}
 \quad(0\leq j\leq t),}
\]
and all other cohomology groups vanish. The complex is zero in $D(R)$ exactly when at least one label is a unit modulo $p$.
\end{theorem}
\begin{proof}
If $h=n$, all entries are zero in $R$. Otherwise choose an entry of minimal valuation, scale its unit part away, and subtract suitable multiples of it from the other entries. These operations form an invertible change of the coefficient row and produce $(p^h,0,\ldots,0)$. The exterior powers of the basis change give a Koszul chain isomorphism, with the convention of \cref{thm:integral-gcd} and \cite{StacksKoszul}.

Each zero entry contributes $R\oplus R[1]$. This proves the complex decomposition. Multiplication by $p^h$ on $R$ has cokernel $Q_h$ and kernel $p^{n-h}R$, the latter isomorphic to $Q_h$ by multiplication by $p^{n-h}$. Hence $C_R(p^h)$ has $Q_h$ in degrees $-1$ and $0$. Counting the shifted copies and using Pascal's identity gives the formula. The unit case is contractible, and otherwise $H^0=Q_h\ne0$.
\end{proof}

For two events the cohomology ranks over $Q_h$ are $1,2,1$ in degrees $-2,-1,0$, rather than the $0,1,1$ pattern over the integral or complete domain. These extra classes are real finite-level cohomology. They may not be discarded while retaining the same complex. In particular, a cohomology group appearing here is not automatically a perfect $R$-module.

\begin{theorem}[Nilpotent interaction forces excisive collapse]\label{thm:nilpotent-collapse}
For any commutative ring $R$ and nilpotent $x\in R$, the thick subcategory generated by $C_R(x)$ is all of $\operatorname{Perf}(R)$. Therefore, over $R=\mathbb Z/p^n$, forcing the interaction of two nonunit scalar labels to vanish in an exact quotient forces that quotient to be zero.
\end{theorem}
\begin{proof}
Choose $N$ with $x^N=0$. For composable scalar multiplications, the octahedral cofiber relation builds $C_R(x^k)$ from $C_R(x^{k-1})$ and $C_R(x)$. Induction puts $C_R(x^N)$ in $\operatorname{Thick}(C_R(x))$. But
\[
 C_R(x^N)=C_R(0)\cong R\oplus R[1].
\]
Retract closure therefore includes $R$, and $R$ thickly generates $\operatorname{Perf}(R)$. This cone-of-a-product mechanism is also part of the standard Koszul calculus \cite{StacksKoszul}.

For two nonunit labels, \cref{thm:prime-power-koszul} makes their interaction a sum of $C_R(p^h)$ and its shift, with $1\leq h\leq n$. Thus it contains $C_R(p^h)$ as a retract. This element is nilpotent. Any exact functor annihilating the interaction annihilates the thick subcategory it generates, hence every perfect object.
\end{proof}

This is stronger than the zero-label counterexample over $\mathbb Z$. At finite $p$-power resolution, even nonzero labels such as $p$ and $p^2$ can force total collapse. Localizing $\mathbb Z/p^n$ at $p$ gives the zero ring. Noncollapse over the integral coefficients does not survive this base change automatically.

\subsection{Several primes give a nonzero, but selective, quotient}

Specify finitely many positive scalar labels and a collection $\mathcal I$ of allowed independent pairs. Include their unit-seed interaction diagrams as the designated tests. This hypothesis ensures that the minimal kernel below is not being inferred from templates that accidentally annihilate an interaction already. Put
\[
 S=\{p\in\mathcal P:\text{some }(a,b)\in\mathcal I
                    \text{ satisfies }p\mid a\text{ and }p\mid b\},
 \qquad D_0=\prod_{p\in\mathcal P\setminus S}p.
\]
Let $\mathcal J_n$ be the tensor-invariant thick closure of those interactions in $\operatorname{Perf}(B_n)$. Positive lengths in \cref{thm:weighted-adic-exchange} do not change $S$.

\begin{theorem}[Prime-component quotient and tower compatibility]\label{thm:finite-prime-quotient}
Under the Chinese remainder equivalence,
\[
 \boxed{\mathcal J_n=
 \prod_{p\in S}\operatorname{Perf}(\mathbb Z/p^n)
 \times\prod_{p\in\mathcal P\setminus S}0.}
\]
In particular there is an equivalence of stable categories
\[
 \operatorname{Perf}(B_n)/\mathcal J_n
 \simeq \prod_{p\in\mathcal P\setminus S}\operatorname{Perf}(\mathbb Z/p^n)
 \simeq\operatorname{Perf}(\mathbb Z/D_0^n).
\]
The equivalence is realized by projection to the retained prime components. Its kernel is constructor-invariant and preserved by every coarsening. If $q_n$ denotes projection, then
\[
 q_nC_{m,n}\simeq C^0_{m,n}q_m
\]
with the supplied base-change coherence. The quotient is nonzero exactly when $S\ne\mathcal P$. When $S=\mathcal P$, use the zero category and the zero ring $\mathbb Z/1$.
\end{theorem}
\begin{proof}
At a prime outside $S$, each designated interaction has a unit label and is acyclic. At a prime in $S$, at least one pair has both labels nonunit. Tensor by the corresponding Chinese remainder idempotent to isolate that component. The nilpotent-collapse theorem then puts its unit in $\mathcal J_n$. Its whole perfect category follows. These two observations prove equality of the kernel with the displayed product subcategory.

Projection to the complementary factors is exact and has an exact fully faithful section. Every object is a direct sum of its selected and retained components. In the quotient the selected summand is zero, so that object becomes isomorphic to its retained summand. Projection and section therefore induce inverse equivalences. This proves the whole quotient assertion directly, without an unproved identification of an arbitrary Verdier quotient with localized perfect complexes.

Tensor products and reduction respect the Chinese remainder idempotents. Projection consequently commutes with reduction, and the same tensor associativities give coherence for composite reductions. All prescribed square interactions vanish after projection, so the earlier excisive reconstruction results apply at every level.
\end{proof}

This repair does not retain the selected prime's finer torsion while deleting only one common exponent. It removes the entire selected prime component. Thus its minimality is a statement about exact excision, not a guarantee that it preserves all arithmetic information desired by a later application.


\section{Completion defects and arithmetic reconstruction}\label{sec:completion-reconstruction}

\subsection{A controlled reconstruction of a resolution tower}

Write
\[
 \widehat B=\varprojlim_n B_n\cong\prod_{p\in\mathcal P}\mathbb Z_p.
\]
The isomorphism follows by taking inverse limits of the compatible Chinese remainder decompositions. Here $\mathbb Z_p$ is the ring of compatible residue sequences modulo $p^n$, and $\mathbb Q_p=\mathbb Z_p[1/p]$.

\begin{theorem}[Uniform based-complex tower reconstruction]\label{thm:based-tower-reconstruction}
Suppose $M_n$ are based finite free cochain complexes over $B_n$, concentrated in a common finite interval with fixed ranks in each degree. Suppose their differential matrices reduce exactly from level $m$ to level $n$ for $m\geq n$. Then there is a unique based complex $\widehat M$ with those ranks over $\widehat B$ whose matrices reduce to the supplied ones.

Compatible chain maps and specified compatible chain homotopies reconstruct coefficientwise. Moreover, as complexes of abelian groups,
\[
 \widehat M\simeq R\varprojlim_n M_n.
\]
Finite tensor products, finite direct sums, shifts, and cones of compatible based maps commute with this reconstruction. In particular this holds for the scalar residual and interaction complexes with fixed integral labels.
\end{theorem}
\begin{proof}
Every matrix entry is a compatible residue sequence and therefore determines a unique element of $\widehat B$. The equation $d^{i+1}d^i=0$ holds because it holds modulo every $D^n$ and the inverse-limit ring is separated. The same argument proves the equations for a reconstructed chain map or homotopy.

Degreewise transition maps of the finite free modules are surjective. In the product model of the homotopy inverse limit, the map
\[
 1-\mathrm{shift}:\prod_n M_n^i\longrightarrow\prod_n M_n^i
\]
is surjective in each degree: choose an initial entry and successively lift the required difference across the surjective transition map. Its kernel is the ordinary inverse limit. Thus the homotopy fiber is represented by that kernel complex. This is the degreewise Mittag--Leffler case of derived inverse limits \cite{StacksRlim}.

All the listed finite constructions have matrices formed by finite sums and products, or by fixed block arrangements of entries. These operations commute with the coefficient inverse limit. This proves the final assertions.
\end{proof}

This theorem does not assert that arbitrary objects chosen independently at every level have a compatible system of bases, bounded ranks, or a coherent lift. It also does not claim an equivalence between all categories of perfect complexes and an inverse limit of categories. Those stronger effectivity statements require additional work. The precise input here is a compatible, uniformly finite, based tower.

\begin{proposition}[Transient negative cohomology in the resolution direction]\label{prop:transient-koszul}
Fix $p$ and an integer $a\geq1$. For $n\geq a$ put $R_n=\mathbb Z/p^n$. Then
\[
 H^0C_{R_n}(p^a)\cong\mathbb Z/p^a,
 \qquad H^{-1}C_{R_n}(p^a)\cong\mathbb Z/p^a.
\]
Under reduction from $m$ to $n$ with $m\geq n\geq a$, the map on $H^0$ is the identity in these coordinates. On $H^{-1}$ it is multiplication by $p^{m-n}$. The negative-cohomology system is therefore pro-zero, although every one of its terms is nonzero. The completed complex is
\[
 C_{\mathbb Z_p}(p^a),\qquad
 H^{-1}=0,\quad H^0=\mathbb Z/p^a.
\]
\end{proposition}
\begin{proof}
Represent $H^{-1}$ by the annihilator $p^{n-a}R_n$ and identify $\mathbb Z/p^a$ with it by $x\mapsto p^{n-a}x$. Reduction of $p^{m-a}x$ gives $p^{n-a}p^{m-n}x$. Thus for each target level $n$, every source level $m\geq n+a$ maps to zero. This is exactly the pro-zero assertion. The cokernel coordinates reduce identically. The preceding theorem reconstructs the complex coefficientwise. Multiplication by $p^a$ on $\mathbb Z_p$ is injective, so its completed kernel vanishes.
\end{proof}

This calculation concerns the actual transition maps on homology, not merely the group appearing at each level. A list of isomorphism types $(\mathbb Z/p^a,\mathbb Z/p^a,\ldots)$ cannot distinguish the identity system from the pro-zero system. The two-axis data must retain those maps.

\subsection{Completion and excisive repair do not generally commute}

Let $S\subseteq\mathcal P$ be the interaction-prime set in \cref{thm:finite-prime-quotient}, and set $T=\prod_{p\in S}p$, with $T=1$ if $S$ is empty. Localizing the finite rings and then completing gives
\[
 \varprojlim_n B_n[1/T]\cong\prod_{p\in\mathcal P\setminus S}\mathbb Z_p.
\]
Completing and then localizing instead gives
\[
 \widehat B[1/T]\cong
 \prod_{p\in S}\mathbb Q_p\times
 \prod_{p\in\mathcal P\setminus S}\mathbb Z_p.
\]
The latter is also an exact realization that annihilates the specified scalar interactions, because each selected prime becomes invertible. There is a canonical comparison map between these rings induced by the reductions. Localization is the usual algebraic operation \cite{StacksLocalization}, not an additional choice of quotient semantics.

\begin{theorem}[Explicit completion--repair defect]\label{thm:completion-repair-defect}
The canonical map
\[
 \chi:\widehat B[1/T]\longrightarrow\varprojlim_n B_n[1/T]
\]
is the projection that forgets the factors $\mathbb Q_p$ for $p\in S$. As $\widehat B$-modules it fits into the split exact sequence
\[
 0\longrightarrow\prod_{p\in S}\mathbb Q_p
 \longrightarrow\widehat B[1/T]
 \xrightarrow{\chi}\prod_{p\in\mathcal P\setminus S}\mathbb Z_p
 \longrightarrow0.
\]
The splitting is a module splitting, not a claimed unital ring section.

For every perfect $\widehat B$-complex $M$, the natural comparison
\[
 \chi_M:M\otimes^{\mathbf L}_{\widehat B}\widehat B[1/T]
 \longrightarrow R\varprojlim_n
 \bigl(M\otimes^{\mathbf L}_{\widehat B}B_n[1/T]\bigr)
\]
has fiber
\[
 \boxed{\operatorname{Fib}(\chi_M)\simeq
 \bigoplus_{p\in S}M\otimes^{\mathbf L}_{\widehat B}\mathbb Q_p.}
\]
Thus $\chi_M$ is an equivalence exactly when every displayed generic fiber is zero. Its cofiber is the shift by $[1]$ of that fiber. These statements take place in $D(\widehat B)$; the field summands need not be perfect over $\widehat B$.
\end{theorem}
\begin{proof}
In the $p$-primary factor of $B_n$, inverting $T$ gives zero if $p\in S$ and does nothing otherwise. In the $p$-primary factor of $\widehat B$, it gives $\mathbb Q_p$ if $p\in S$ and leaves $\mathbb Z_p$ otherwise. This proves both ring formulas. The induced map is zero on the first set of factors and the identity on the others, proving the split sequence.

For $M=\widehat B$, the derived inverse limit on the right is the ordinary one, since the retained residue-module tower is degreewise surjective. The asserted comparison and fiber are therefore exactly the split sequence. For a general perfect $M$, all functors in the claimed identity are exact and preserve finite sums and retracts. Homotopy inverse limits are exact in a stable derived category. The natural comparison identifies its right side with the retained summand of $M$: this holds on the unit by the just-computed inverse limit and hence on its thick closure, which is $\operatorname{Perf}(\widehat B)$. Tensoring the split sequence now identifies the fiber with the indicated finite sum. A morphism in a stable category is an equivalence precisely when its fiber is zero.
\end{proof}

\begin{example}[A nonzero generic fiber is lost at every finite repair]\label{ex:repair-before-completion}
For $\mathcal P=\{p\}$ and two independent labels divisible by $p$, every finite excisive quotient is zero, whereas
\[
 \bigl(\varprojlim_n\mathbb Z/p^n\bigr)[1/p]=\mathbb Q_p\ne0.
\]
For a nonzero finite-level quotient, take $D=30$ and labels $6,10$. Only the prime $2$ is shared, so the finite repair is
\[
 B_n[1/2]\cong\mathbb Z/15^n.
\]
The two completed outputs are $\mathbb Q_2\times\mathbb Z_3\times\mathbb Z_5$ and $\mathbb Z_3\times\mathbb Z_5$, with fiber $\mathbb Q_2$. This is an explicit failure of an unconditional bi-completion principle, even though the finite quotient commutes with every finite coarsening.
\end{example}

The generic-fiber criterion also gives a positive case. For $M=C_{\widehat B}(u)$ with a nonzero integer $u$, its selected generic fibers vanish, because $u$ becomes invertible in each $\mathbb Q_p$. Thus this comparison is an equivalence on that torsion residual even though it is not an equivalence on the unit. Commutation depends on the object, not only on the existence of coherent finite-level squares. The distinction is consistent with the standard derived-completion fact that a complex on which a defining element is invertible has zero completion \cite{StacksDerivedCompletion}.

\subsection{Reconstruction that retains the overlap instead of deleting it}

There is a different operation from killing an interaction: retain the localized and completed descriptions together with their common generic information. In the present arithmetic setting this yields an actual pullback theorem.

Put
\[
 A=\mathbb Z[1/D],\qquad B=\prod_{p\in\mathcal P}\mathbb Z_p,
 \qquad G=\prod_{p\in\mathcal P}\mathbb Q_p.
\]

\begin{theorem}[Arithmetic fracture reconstruction]\label{thm:arithmetic-fracture}
The diagonal maps and their difference give an exact sequence of abelian groups
\[
 0\longrightarrow\mathbb Z
 \longrightarrow A\oplus B
 \xrightarrow{(a,b)\mapsto a-b}G
 \longrightarrow0.
\]
Consequently, for every perfect integral complex $M$ there is a natural homotopy pullback square in $D(\mathbb Z)$,
\[
 \begin{tikzcd}
 M\arrow[r]\arrow[d]&M\otimes^{\mathbf L}_{\mathbb Z}A\arrow[d]\\
 M\otimes^{\mathbf L}_{\mathbb Z}B\arrow[r]&M\otimes^{\mathbf L}_{\mathbb Z}G.
 \end{tikzcd}
\]
Its bottom-left term is $R\varprojlim_n(M\otimes^{\mathbf L}_{\mathbb Z}B_n)$. In particular $M$ is zero if and only if its localized and completed terms are both zero.
\end{theorem}
\begin{proof}
An element of $A$ integral at every $p\in\mathcal P$ is an integer: its possible denominator primes all lie in $\mathcal P$, and integrality removes each of them. This identifies the kernel with the diagonal copy of $\mathbb Z$.

For surjectivity take $(y_p)\in G$. Choose $N$ large enough that $D^Ny_p\in\mathbb Z_p$ for every $p$. For each $p$, choose the residue of $D^Ny_p$ modulo $p^N$. The ordinary Chinese remainder theorem provides $z\in\mathbb Z$ with all these residues. Set $a=z/D^N\in A$. Then $a-y_p\in\mathbb Z_p$ for every $p$, so $(a,(a-y_p)_p)$ maps to $(y_p)_p$. This proves exactness.

The exact sequence gives a cofiber sequence in the derived category. Derived tensor with $M$ is exact, so it gives the displayed homotopy pullback. For a bounded finite free representative of $M$, the tower of reductions has uniform ranks and compatible matrices. The based-tower theorem identifies its derived limit with $M\otimes B$. Finally if the localized and completed terms are zero, so is their generic base change, and the pullback is zero. The converse is immediate.
\end{proof}

This is the classical arithmetic-square mechanism, proved directly here using the Chinese remainder theorem and derived tensor. It is not claimed as a new local--global correspondence. It does show exactly which overlap data must accompany these arithmetic representations to reconstruct a supplied perfect complex.

\begin{corollary}[Residuals reconstruct through the same arithmetic square]\label{cor:residual-fracture}
Let $X$ be a finite coherent event diagram in $\operatorname{Perf}(\mathbb Z)$. Every edge residual, square interaction, and initial-to-terminal residual of $X$ satisfies \cref{thm:arithmetic-fracture}. Its terms are the corresponding residuals of the diagrams obtained by base change to $A,B,G$. Each such integral residual vanishes if and only if both its localized and completed versions vanish.
\end{corollary}
\begin{proof}
Derived base change is exact and therefore preserves all the finite cofibers defining those residuals. Apply the arithmetic fracture theorem to the resulting perfect integral complex. Naturality gives the maps on the whole residual diagram, not merely equality of its cohomology dimensions.
\end{proof}

\subsection{Effective gluing of lattices from independent local input}

The previous theorem starts with a global complex. The following finite-rank theorem instead constructs a global object from supplied local and localized data.

\begin{theorem}[Effective arithmetic lattice gluing]\label{thm:lattice-gluing}
Let $V$ be a finite free $A=\mathbb Z[1/D]$-module of rank $r$. For each $p\in\mathcal P$ let $L_p$ be a finite free $\mathbb Z_p$-module, with an isomorphism
\[
 \theta_p:V\otimes_A\mathbb Q_p\xrightarrow{\sim}
 L_p\otimes_{\mathbb Z_p}\mathbb Q_p.
\]
Then
\[
 \Lambda=\{v\in V:\theta_p(v)\in L_p\text{ for every }p\in\mathcal P\}
\]
is a finite free $\mathbb Z$-module of rank $r$, and the natural maps identify
\[
 \Lambda[1/D]\cong V,\qquad
 \Lambda\otimes_{\mathbb Z}\mathbb Z_p\cong L_p.
\]
It is the unique lattice inside $V$ with these properties and these specified embeddings. Compatible maps of gluing data restrict uniquely to maps of lattices. Hence finite free integral modules are equivalent to this category of localized free modules, local lattices, and specified generic identifications.
\end{theorem}
\begin{proof}
Choose a basis $V=A^r$, and let $E_p=\theta_p^{-1}(L_p)$ be the corresponding lattice in $\mathbb Q_p^r$. Bounds on the entries of $\theta_p$ and its inverse give a common $N\geq1$ such that, with $b=D^N$,
\[
 b\mathbb Z_p^r\subseteq E_p\subseteq b^{-1}\mathbb Z_p^r
 \quad\text{for all }p.
\]
The intersection defining $\Lambda$ therefore satisfies
\[
 b\mathbb Z^r\subseteq\Lambda\subseteq b^{-1}\mathbb Z^r.
\]
The right inclusion uses the same integrality-of-denominators argument as the arithmetic fracture theorem, coordinatewise. Consequently $\Lambda$ is a finite-rank free abelian lattice of rank $r$, and localizing the bounds at $D$ gives $\Lambda[1/D]=A^r$.

Here is also an explicit finite construction proving the local identifications. Put $F_p=bE_p$, so
\[
 b^2\mathbb Z_p^r\subseteq F_p\subseteq\mathbb Z_p^r.
\]
The Chinese remainder isomorphism identifies
\[
 \mathbb Z^r/b^2\mathbb Z^r\cong
 \prod_{p\in\mathcal P}\mathbb Z_p^r/b^2\mathbb Z_p^r.
\]
Let $H\subseteq\mathbb Z^r$ be the inverse image of the product of the subgroups $F_p/b^2\mathbb Z_p^r$. Directly from the definitions, $H=b\Lambda$. The quotient $\mathbb Z^r/H$ is the product of the corresponding finite $p$-primary quotients. Tensoring with $\mathbb Z_p$ retains exactly the $p$ component. Thus $H\otimes\mathbb Z_p$ is the inverse image of $F_p/b^2\mathbb Z_p^r$, namely $F_p$. Dividing by $b$ gives $\Lambda\otimes\mathbb Z_p=E_p$, and $\theta_p$ identifies it with $L_p$.

For completeness, tensoring the displayed subgroup inclusion is exact: $\mathbb Z_p$ is torsion-free over $\mathbb Z$, so it is flat; alternatively the assertion follows from a Smith diagonal presentation of the finite-index subgroup $H$. The latter proves the same statement using only integer bases.

Any other lattice with the specified generic and local identifications has analogous bounds by powers of $D$. After a common scaling, its finite quotient is determined by the same primary subgroups in the Chinese remainder decomposition. It therefore equals $\Lambda$. A compatible $A$-linear map $V\to V'$ sends every $\theta_p$-integral element to a $\theta'_p$-integral element and hence restricts to $\Lambda\to\Lambda'$. Conversely a lattice map induces the compatible generic and local maps. Uniqueness follows from the embedding in $V$. These constructions are inverse on objects up to the specified identifications and on morphisms, proving the equivalence. The zero-rank case gives the zero module throughout.
\end{proof}

This is an effective local--global existence theorem for lattices, not for Galois representations or automorphic representations. The generic identifications are essential input. The construction does not derive them from matching coarse invariants. Nor does it assert gluing of arbitrary nonlinear structures, group actions, filtrations, or analytic conditions. Those would have to descend through the lattice construction with their own compatibility proofs.

\paragraph{Consequences for the calculus.}
The new example separates three operations that a universal slogan would conflate. Exact finite-level excision can erase an entire prime component. Completion can retain a generic fiber that all those finite repairs erased. Effective reconstruction instead retains the localized object, the completed objects, and their overlap identifications. The lattice theorem shows that this latter pattern can genuinely construct a global arithmetic object from local input, in a precisely specified class.

\begin{corollary}[Compatible operators glue with the lattice]\label{cor:lattice-operators}
In \cref{thm:lattice-gluing}, suppose $f_1,\ldots,f_s\in\operatorname{End}_A(V)$ and $g_{p,1},\ldots,g_{p,s}\in\operatorname{End}_{\mathbb Z_p}(L_p)$ satisfy the exact intertwining equations
\[
 \theta_p(f_i\otimes1)=(g_{p,i}\otimes1)\theta_p.
\]
Then they induce unique integral endomorphisms of $\Lambda$. Every specified noncommutative polynomial relation with integer coefficients holding on $V$ holds on $\Lambda$. If the operators and their inverses are supplied compatibly, they glue as automorphisms.
\end{corollary}
\begin{proof}
For $v\in\Lambda$, the displayed identity puts $\theta_p(f_i v)=g_{p,i}\theta_p(v)$ in $L_p$ for every $p$. Hence $f_i v\in\Lambda$. Restriction is unique because $\Lambda$ spans $V$ after localization. The polynomial relations hold by restriction from $V$, and applying the argument to the supplied inverses gives automorphisms.
\end{proof}

This corollary glues an already compatible operator action. It does not construct a generic action, a continuous Galois action, or an automorphic realization from scalar or Frobenius invariants alone.


\section{Local-to-global reconstruction}

\subsection{A \v{C}ech torsor residual as a nonstable obstruction test bed}

The stable Verdier quotient is not the only setting in which the residual/repair/obstruction pattern appears. A classical torsor gives a minimal nonstable model.

\begin{theorem}[Torsor residual obstruction]\label{thm:torsor-obstruction}
Let $\mathcal A$ be a sheaf of abelian groups on a space or a site with the fiber products used below, let $P$ be an $\mathcal A$-torsor, and let $\{U_i\to U\}$ be a cover on which $P$ admits local sections $s_i\in P(U_i)$. On overlaps there are unique elements
\[
a_{ij}\in\mathcal A(U_{ij})
\]
such that
\[
s_j=s_i+a_{ij}.
\]
Then:
\begin{enumerate}
    \item $a=(a_{ij})$ is a \v{C}ech $1$-cocycle;
    \item changing the local sections by a $0$-cochain $b=(b_i)$,
    \[
    s_i'=s_i+b_i,
    \]
    changes the residual by the \v{C}ech coboundary
    \[
    a'_{ij}=a_{ij}+b_j-b_i;
    \]
    \item the class
    \[
    \Omega(P;\{U_i\})=[a]\in\check H^1(\{U_i\},\mathcal A)
    \]
    is independent of the chosen local sections;
    \item $\Omega(P;\{U_i\})=0$ if and only if the local sections can be modified by admissible local repairs so that they glue to a global section of $P$ over $U$.
\end{enumerate}
\end{theorem}

\begin{proof}
On a triple overlap,
\[
s_k=s_j+a_{jk}=s_i+a_{ij}+a_{jk},
\]
while also $s_k=s_i+a_{ik}$. Freeness of the torsor action gives
\[
a_{ij}+a_{jk}=a_{ik},
\]
which is the \v{C}ech cocycle identity.

If $s_i'=s_i+b_i$, then
\[
s_j'=s_j+b_j=s_i+a_{ij}+b_j
=s_i'+(a_{ij}+b_j-b_i),
\]
so the residual changes by the coboundary $\delta b$. Hence its cohomology class is independent of the chosen local sections.

If $[a]=0$, choose $b$ so that $a+\delta b=0$. The repaired sections $s_i'=s_i+b_i$ then agree on every overlap, hence glue by the sheaf property to a global section of $P$. Conversely, if a global section exists, its restrictions have zero transition residual. Comparing those restrictions with the original $s_i$ shows that $a$ differs from zero by a coboundary, so $[a]=0$.
\end{proof}

\status{proved directly from torsor descent and \v{C}ech cohomology}

For the standard torsor interpretation see \cite{StacksTorsors}.

In Komorebi language,
\[
\boxed{
\text{overlap residual}
\;/\;
\text{local repair}
\;=\;
\text{global obstruction class}.}
\]
Here the residual is $a_{ij}$, admissible repair is the change $s_i\mapsto s_i+b_i$, repairable residual cocycles are coboundaries, and the obstruction is the cohomology class $[a]$. This local repair relation is not being identified with the monoid of exact endofunctors in the stable quotient construction. This example is important because it realizes the proposed obstruction pattern before stabilization and shows that the abstract vocabulary recovers an established local-to-global obstruction mechanism rather than manufacturing one artificially.
It is specifically an obstruction to reconstructing a global \emph{section} (equivalently, to trivializing the torsor), not an obstruction to the existence of the torsor itself.

\subsection{Descent data}

Let $(\mathcal C,J)$ be a site and let $\mathscr K$ assign a Komorebi category to each object of the site, with restriction functors on overlaps.

Given a cover $\{U_i\to U\}$, local objects $X_i\in\mathscr K(U_i)$ carry descent isomorphisms
\[
\phi_{ij}:X_i|_{U_{ij}}\xrightarrow{\sim}X_j|_{U_{ij}}
\]
satisfying the cocycle condition on triple overlaps.

\begin{definition}[Effective descent]
Fix a cover and its complete \v{C}ech diagram, including all iterated overlaps. Effective descent means that the natural functor from the global category to the category of descent data is an equivalence. Essential surjectivity supplies global objects. Full faithfulness supplies descent of morphisms. Uniqueness of object classes alone is insufficient.
\end{definition}

\begin{theorem}[Descent transfer lemma]\label{thm:descent-transfer}
Let $A^\bullet,B^\bullet:\Delta\to\Cat_\infty$ be small \v{C}ech descent diagrams for a fixed cover. Suppose $\Lambda^\bullet:A^\bullet\to B^\bullet$ is a natural transformation that is an equivalence at every degree, including every overlap level. Then
\[
\lim_{[n]\in\Delta} A^n\simeq\lim_{[n]\in\Delta} B^n.
\]
For ordinary category-valued descent, use the corresponding coherent pseudolimits rather than strict equalizers of object sets.
\end{theorem}

\begin{proof}
Equivalences in a functor category are detected objectwise. Hence $\Lambda^\bullet$ is an equivalence of diagrams. Applying the limit functor preserves this equivalence. The ordinary category formulation is the same argument with coherent descent categories and their quasi-inverses.
\end{proof}

\begin{corollary}[Global reconstruction transfer]\label{cor:global-transfer}
If the maps from the two global categories to these descent categories are equivalences, then the induced equivalence of descent categories yields an equivalence of global categories.
\end{corollary}

This is the abstract form of the principle:
\[
\boxed{
\text{local equivalence}
+
\text{restriction compatibility}
+
\text{effective descent}
\Longrightarrow
\text{global equivalence}.
}
\]

\section{Seeded local-global reciprocity blueprint}

The previous section suggests a precise architecture for Langlands-facing work without claiming that the arithmetic hypotheses have been discharged.

Let
\[
\mathbb K_{\mathrm{arith}}
\qquad\text{and}\qquad
\mathbb K_{\mathrm{aut}}
\]
be arithmetic and automorphic seeded two-axis calculi.

A strong correspondence should include:
\begin{align*}
\Lambda(TX)&\simeq T'\Lambda(X),\\
\Lambda(CX)&\simeq C'\Lambda(X),\\
\Lambda(\Delta_{T,C}X)&\simeq\Delta_{T',C'}\Lambda(X),\\
\bar\Lambda(\Obs X)&\simeq\Obs'\Lambda(X).
\end{align*}

Here $\bar\Lambda$ is the induced functor on the designated obstruction quotients. The last two requirements are stronger than matching ordinary local invariants. They demand that the correspondence preserve the structure of failed exchange and failed reconstruction.

\begin{conjecture}[Seeded reciprocity principle]\label{conj:seeded-reciprocity}
Suppose arithmetic and automorphic Komorebi systems satisfy:
\begin{enumerate}
    \item local equivalence at every required place and resolution;
    \item compatibility with resolution restriction;
    \item preservation of the relevant local invariants;
    \item transport of mixed residual and obstruction data;
    \item effective descent on both sides;
    \item a global separation principle ensuring uniqueness from sufficiently rich local data.
\end{enumerate}
Then the induced local equivalence extends to an equivalence of the corresponding admissible global categories.
\end{conjecture}

Under categorical formulations strong enough to make items (1)--(5) an equivalence of descent systems, the conclusion follows from \cref{cor:global-transfer}. The research content lies in proving that the arithmetic and automorphic systems really instantiate those hypotheses.

\paragraph{Interpretation.}
This reframes a possible route to local-global reciprocity as:
\[
\boxed{
\text{identify the globalization obstruction}
\quad\text{and prove that local correspondence transports it.}
}
\]
If
\[
\Lambda_*(\Omega_{\mathrm{arith}})=\Omega_{\mathrm{aut}},
\]
and $\Lambda_*$ preserves and reflects the zero object, then vanishing of one obstruction is equivalent to vanishing of the other. An equivalence, or a conservative exact functor, suffices for this extra detection condition. A general exact functor only preserves vanishing. Effective reconstruction then gives the stated global conclusion. Neither a sheaf/stack model for arbitrary local Langlands data nor automorphic existence follows from the abstract categorical lemmas. Places of a number field are not automatically a covering family for an effective descent problem. Those arithmetic identifications remain obligations.

\section{Toy models}

\subsection{\texorpdfstring{$p$}{p}-adic congruence resolution}

Let the resolution level be $n\ge1$ and take
\[
X_n=\mathbb Z/p^n\mathbb Z.
\]
For $m\le n$, coarsening is reduction
\[
C_{n\to m}:\mathbb Z/p^n\mathbb Z\to\mathbb Z/p^m\mathbb Z.
\]
For any polynomial $f\in\mathbb Z[x]$, the induced transformation satisfies
\[
C_{n\to m}(f(x))
=
f(C_{n\to m}(x)).
\]
Thus the transformation-resolution square commutes exactly. This is a set-level equality, not a cofiber calculation. Regarding the sets as discrete categories gives an identity comparison, while a stable residual requires a separately specified stable realization.
This is a simple instance where arithmetic transformation is resolution-natural.

By contrast, operations involving division by $p$ need not descend uniformly. Their domain of definition and lift multiplicity depend on valuation. This gives a natural source of nontrivial resolution obstruction.

\subsection{Dyadic averaging}

Let a fine state be
\[
x=(x_1,x_2)\in\mathbb R^2
\]
and let coarse resolution record only the average
\[
C(x)=\frac{x_1+x_2}{2}.
\]
Let the fine transformation be
\[
A_1(x_1,x_2)=(ax_1,bx_2)
\]
and the coarse transformation be multiplication by $c$:
\[
A_0(y)=cy.
\]
Then
\[
C(A_1x)-A_0(Cx)
=
\frac{(a-c)x_1+(b-c)x_2}{2}.
\]
Choosing
\[
c=\frac{a+b}{2}
\]
gives
\[
C(A_1x)-A_0(Cx)
=
\frac{a-b}{4}(x_1-x_2).
\]
The defect depends only on the antisymmetric fine mode $x_1-x_2$, precisely the information erased by averaging. This is a concrete model of a mixed residual exposing structure invisible at coarse resolution.

\subsection{Exact symmetry under coarsening}

For dyadic vectors, pairwise averaging commutes with reversal up to reversal of the coarse vector. This gives an exact exchange law and hence a zero mixed residual. The contrast with the previous example illustrates that residuals depend jointly on the transformation and the resolution operator.

\section{Derived invariants for the zoo}

The two-axis formalization suggests new quantities that can be defined without changing the underlying category theory.

\subsection{Commutation depth}

Fix $X\in\Tcat_s$. First record the exact commutation profile
\[
\mathcal E_A(X;s)=\{r\le s:\Delta^A_{s,r}(X)\simeq0\}.
\]
No interval or monotonicity property is assumed. For $R=\mathbb N$ and finite $s$, a well-defined consecutive coarsening depth is
\[
\kappa_A(X;s):=\max\{k\in\{0,\ldots,s\}:\Delta^A_{s,r}(X)\simeq0
\text{ for all }s-k\le r\le s\}.
\]
The set is nonempty because $\Delta^A_{s,s}(X)\simeq0$. A larger value means that more consecutive coarsenings preserve this exchange. The preceding version's supremum over exact target levels was uninformative on a fixed finite tower because the identity level $r=s$ is always exact.

\subsection{Residual complexity}

When a numerical size function $\|\cdot\|$ is available on the fiber, define
\[
\mathcal E_A(s,r;X):=\|\Delta^A_{s,r}(X)\|.
\]
Possible choices include dimension, rank, norm, entropy, categorical amplitude, support size, or proof-complexity measures.

\subsection{Resolution branching profile}

For a chain $r_0<r_1<\cdots<r_n$ and chosen lifts $X_i\in\Tcat_{r_i}$ with $C_{r_{i+1}\to r_i}X_{i+1}\cong X_i$, record
\[
B_X=(b_{r_0\to r_1}(X_0),\ldots,b_{r_{n-1}\to r_n}(X_{n-1})).
\]
This profile measures how hidden multiplicity emerges under refinement.

\subsection{Zoo growth spectrum}

Instead of a single growth function, separate growth by generator type:
\[
g_T(n),\qquad g_R(n),\qquad g_\Delta(n).
\]
This may distinguish theories in which most complexity comes from representation changes, resolution branching, or residual generation.

\section{Candidate axioms for a foundational version}

The present note derives its first results from ordinary category and order theory. A stronger foundational program would isolate Komorebi-specific axioms.

\begin{axiom}[Two-axis presentation]
Every admissible structured state is represented in a system
\[
\pi:\K\to\Rcat
\]
whose fibers encode transformations and whose base encodes resolution.
\end{axiom}

\begin{axiom}[Resolution compositionality]
Coarsening is functorial:
\[
C_{t\to r}=C_{s\to r}C_{t\to s}.
\]
\end{axiom}

\begin{axiom}[Exchange data]
Every designated transformation family carries a coherent comparison
\[
A_rC_{s\to r}\Longrightarrow C_{s\to r}A_s.
\]
\end{axiom}

\begin{axiom}[Residual closure]
Every noninvertible exchange comparison generates an admissible residual object that may enter the seed zoo.
\end{axiom}

\begin{axiom}[Effective reconstruction, domain-specific]
For a specified class of local systems, vanishing global obstruction is equivalent to effective global reconstruction.
\end{axiom}

These are axioms for a class of structures, not proposed additional axioms of ZFC. Two-axis existence alone imposes no substantive restriction: every category $\mathcal C$ has the degenerate presentation over the one-object resolution category with fiber $\mathcal C$. A foundational claim needs nontrivial, semantically specified resolutions and a theorem relating them to the intended mathematics.

The last axiom should not be imposed universally without justification. In serious applications it should be a theorem of a sheaf, stack, patching, or descent theory.

\section{Machine-native semantics}

The seed calculus can be interpreted as a search grammar. A machine state may contain
\[
\mathsf S=
(\Seeds,\mathscr T,\mathscr R,\Delta,\Omega,\mathsf{Cert}).
\]
For a target problem $P$, search can occur jointly over transformation and resolution:
\[
(T^\ast,R^\ast)
\in
\operatorname*{arg\,min}_{T,R}
\mathcal C\bigl(\Omega(T,R;P)\bigr),
\]
where $\mathcal C$ is a problem-dependent residual or obstruction cost. This expression describes a search objective. It does not assert that a minimizer exists or that the search is decidable.

This changes the computational question from
\[
\text{``Which operation solves the problem?''}
\]
to
\[
\boxed{
\text{``At which resolution and in which representation does the obstruction become simplest?''}
}
\]

A generic discovery loop is:
\[
\boxed{
\text{seed}
\to
\text{generate}
\to
\text{normalize}
\to
\text{measure residual}
\to
\text{promote residual}
\to
\text{expand zoo}
\to
\text{certify}.
}
\]

Formal proof assistants can serve as the certificate layer without requiring the generative search process itself to be trusted.


\section{Universal constructions: syntax and stable residual completion}\label{sec:universal}

This section addresses the first foundational novelty test. The result is deliberately split into a syntactic theorem, which is standard double-category machinery, and a residual-completion theorem, which is a formal consequence of the free stable envelope of an additive category. The combination gives a rigorous universal construction for the additive/stable version of the seeded two-axis calculus. It does \emph{not} by itself establish that the construction is absent from the literature under another name.

\subsection{Free seeded two-axis double syntax}

\begin{definition}[Seeded two-axis double computad]\label{def:twoaxis-computad}
A \emph{seeded two-axis double computad} $\mathbb S$ consists of a small double computad together with:
\begin{enumerate}
    \item a distinguished set $\Seeds\subseteq \mathbb S_0$ of seed objects;
    \item a designated class of horizontal generating $1$-cells, interpreted as transformation generators;
    \item a designated class of vertical generating $1$-cells, interpreted as resolution generators;
    \item designated generating $2$-cells, interpreted as exchange or compatibility cells.
\end{enumerate}
Generator boundaries may be well-typed words, but no composition of the generators is identified with another generator at the signature level. An unlabelled realization of $\mathbb S$ in a strict double category $\mathbb D$ is a map of the underlying double computads. A label-preserving realization additionally requires matching distinguished subsets in the target and preservation of those designations.
\end{definition}

The label data are algebraically passive: they distinguish generators but impose no new equations. Hence the ordinary free-double-category construction applies unchanged.

\begin{theorem}[Free seeded two-axis double category]\label{thm:free-double}
For every small seeded two-axis double computad $\mathbb S$, there exists a strict double category $\mathbb D_{\mathbb S}$ and a canonical map of seeded double computads
\[
\iota_{\mathbb S}:\mathbb S\longrightarrow U\mathbb D_{\mathbb S}
\]
such that for every strict double category $\mathbb D$, restriction along $\iota_{\mathbb S}$ induces a natural bijection
\[
\operatorname{DblCat}(\mathbb D_{\mathbb S},\mathbb D)
\cong
\operatorname{DblCptd}(\mathbb S,U\mathbb D),
\]
for the underlying computadic data. If the target carries distinguished subsets for the same labels, restriction gives the analogous bijection with label-preservation imposed on both sides.
\end{theorem}

\begin{proof}
Form horizontal and vertical words in the generating arrows. Form square terms from the generating squares and identity squares by every well-typed horizontal and vertical composition. Quotient the terms by the congruence generated by units, associativity, and interchange. These relations give a strict double category. Every interpretation of the generating data extends by evaluation of terms, and evaluation respects the congruence because the target satisfies the same axioms. Its values on generators force uniqueness. Labels merely restrict the allowed generator assignments when corresponding target labels have been specified.
\end{proof}

\status{proved by translation to standard free double-category machinery}

This establishes a genuine universal property for the \emph{two-axis syntax}. It is not a novelty claim: free double categories and double computads are established constructions; see Fairbanks--Shulman~\cite{FairbanksShulman2026}.

\subsection{Additive indexed systems}

The residual layer becomes canonical after passing to stable categories. Let $\mathbf{Cat}^{\mathrm{add}}$ denote the $\infty$-category of small additive $\infty$-categories and additive functors, and let $\mathbf{Cat}^{\mathrm{st}}$ denote small stable $\infty$-categories and exact functors. Ordinary additive categories are included by their nerves. The forgetful functor retains the underlying additive $\infty$-category, not merely its homotopy category.

\begin{definition}[Additive two-axis system]
An \emph{additive two-axis system} over a small resolution category $\Rcat$ is a functor
\[
F:\Rcat^{op}\longrightarrow \mathbf{Cat}^{\mathrm{add}}.
\]
Thus every fiber $F(r)$ is additive and every coarsening functor $C_{s\to r}$ is additive.
\end{definition}

For a small additive category $\mathcal A$, write $\operatorname{St}^{\mathrm{add}}(\mathcal A)$ for its free stable envelope. We use the established adjunction
\[
\operatorname{St}^{\mathrm{add}}:
\mathbf{Cat}^{\mathrm{add}}
\rightleftarrows
\mathbf{Cat}^{\mathrm{st}}:U,
\]
characterized by
\[
\operatorname{Fun}^{\mathrm{ex}}
\bigl(\operatorname{St}^{\mathrm{add}}(\mathcal A),\mathcal C\bigr)
\simeq
\operatorname{Fun}^{\oplus}(\mathcal A,U\mathcal C)
\]
for every stable $\mathcal C$; see Saunier--Winges~\cite{SaunierWinges2026}.

\begin{theorem}[Pointwise stable completion]\label{thm:pointwise-stable}
Let $\Rcat$ be small. Postcomposition with $\operatorname{St}^{\mathrm{add}}$ defines a functor
\[
\operatorname{St}^{\mathrm{add}}_{\Rcat}:
\operatorname{Fun}(\Rcat^{op},\mathbf{Cat}^{\mathrm{add}})
\longrightarrow
\operatorname{Fun}(\Rcat^{op},\mathbf{Cat}^{\mathrm{st}}),
\qquad
F\longmapsto \operatorname{St}^{\mathrm{add}}\circ F.
\]
It is left adjoint to pointwise forgetting
\[
U_{\Rcat}:
\operatorname{Fun}(\Rcat^{op},\mathbf{Cat}^{\mathrm{st}})
\longrightarrow
\operatorname{Fun}(\Rcat^{op},\mathbf{Cat}^{\mathrm{add}}).
\]
Equivalently, for every additive two-axis system $F$ and stable two-axis system $G$ there is a natural equivalence
\[
\operatorname{Map}
\bigl(\operatorname{St}^{\mathrm{add}}_{\Rcat}F,G\bigr)
\simeq
\operatorname{Map}
\bigl(F,U_{\Rcat}G\bigr).
\]
\end{theorem}

\begin{proof}
Any adjunction $L\dashv U$ induces an adjunction on functor categories by postcomposition. Concretely, a natural transformation
\[
\alpha:L\circ F\Rightarrow G
\]
has pointwise adjoint mates
\[
\alpha_r^{\flat}:F(r)\longrightarrow U(G(r)).
\]
Naturality of the unit and counit implies that the family $\{\alpha_r^{\flat}\}_r$ is natural in $r$. The mate construction is invertible and respects higher coherence. Applying this to $L=\operatorname{St}^{\mathrm{add}}$ gives the result.
\end{proof}

\status{proved}

This is the precise sense in which the stable version of the two-axis system has a \emph{free residual-complete envelope}: no extra choices are required to create finite cofibers in the fibers, and every additive map into a stable two-axis system extends essentially uniquely to an exact one.

\subsection{Exchange data and canonical residuals}

\begin{definition}[Coherent additive transformation family]
Let $F:\Rcat^{op}\to\mathbf{Cat}^{\mathrm{add}}$. A coherent additive transformation family $A$ consists of additive endofunctors
\[
A_r:F(r)\to F(r)
\]
and, for each $r\le s$, a natural transformation
\[
\eta^A_{s,r}:A_rC_{s\to r}\Longrightarrow C_{s\to r}A_s,
\]
satisfying the identity and transitivity coherences for chains $r\le s\le t$.
\end{definition}

The universal property above is an equivalence of functor $\infty$-categories, not only of their cores. This distinction is essential for noninvertible exchange maps. To extend $\eta:f\Rightarrow g$ between additive functors $\mathcal A\to\mathcal B$, compose it with $j_{\mathcal B}:\mathcal B\to\operatorname{St}^{\mathrm{add}}(\mathcal B)$ and lift the resulting morphism through
\[
\operatorname{Fun}^{\mathrm{ex}}(\operatorname{St}^{\mathrm{add}}(\mathcal A),\operatorname{St}^{\mathrm{add}}(\mathcal B))
\simeq\operatorname{Fun}^{\oplus}(\mathcal A,\operatorname{St}^{\mathrm{add}}(\mathcal B)).
\]
Full faithfulness transports the supplied composition coherences as well. Accordingly $A_r$, $C_{s\to r}$, and their coherently supplied exchange cells extend, essentially uniquely relative to the input, to exact functors and natural transformations
\[
\widehat A_r,
\qquad
\widehat C_{s\to r},
\qquad
\widehat\eta^A_{s,r}:
\widehat A_r\widehat C_{s\to r}
\Longrightarrow
\widehat C_{s\to r}\widehat A_s.
\]

\begin{definition}[Universal stable mixed residual]\label{def:universal-residual}
For $X\in\operatorname{St}^{\mathrm{add}}(F(s))$, define
\[
\Delta^A_{s,r}(X)
:=
\Cofib\!\left(
\widehat\eta^A_{s,r,X}
\right)
\in
\operatorname{St}^{\mathrm{add}}(F(r)).
\]
\end{definition}

\begin{theorem}[Universal residual preservation]\label{thm:residual-preservation}
Let $G:\Rcat^{op}\to\mathbf{Cat}^{\mathrm{st}}$ be a stable two-axis system and let
\[
\Phi:\operatorname{St}^{\mathrm{add}}_{\Rcat}F\longrightarrow G
\]
be an exact natural transformation which intertwines the chosen transformation families and exchange cells. Then for every $r\le s$ and every $X$ there is a canonical equivalence
\[
\Phi_r\bigl(\Delta^A_{s,r}(X)\bigr)
\simeq
\Cofib\!\left(
\Phi_r(\widehat\eta^A_{s,r,X})
\right).
\]
In particular, if the exchange data in $G$ are denoted by $\eta^B$, then
\[
\Phi_r\bigl(\Delta^A_{s,r}(X)\bigr)
\simeq
\Delta^B_{s,r}(\Phi_sX).
\]
\end{theorem}

\begin{proof}
Each $\Phi_r$ is exact, hence preserves finite colimits and therefore cofiber sequences. Apply $\Phi_r$ to the defining cofiber sequence
\[
\widehat A_r\widehat C_{s\to r}X
\xrightarrow{\widehat\eta^A_{s,r,X}}
\widehat C_{s\to r}\widehat A_sX
\longrightarrow
\Delta^A_{s,r}(X).
\]
The result is a cofiber sequence in $G(r)$. Compatibility of $\Phi$ with the transformation and exchange data identifies its first morphism with the target exchange morphism, giving the stated equivalence.
\end{proof}

\status{proved}

Thus residual formation is not merely available in the stable completion: it is \emph{universal under exact realization}.

\subsection{Residual promotion is forced by stable closure}

\begin{definition}[Stable seeded zoo]
Fix seed sets $\Seeds_r\subseteq\Ob(F(r))$. A \emph{stable seeded sub-zoo} of $\operatorname{St}^{\mathrm{add}}_{\Rcat}F$ is a family of full stable subcategories
\[
\mathcal Z_r\subseteq\operatorname{St}^{\mathrm{add}}(F(r))
\]
such that:
\begin{enumerate}
    \item $\mathcal Z_r$ contains the images of all seeds in $\Seeds_r$;
    \item $\widehat A_r(\mathcal Z_r)\subseteq\mathcal Z_r$ for every designated transformation generator;
    \item $\widehat C_{s\to r}(\mathcal Z_s)\subseteq\mathcal Z_r$ for every $r\le s$.
\end{enumerate}
The \emph{minimal stable zoo} is the intersection of all such families.
\end{definition}

\begin{proposition}[Existence of the minimal stable zoo]\label{prop:min-stable-zoo}
The minimal stable zoo exists.
\end{proposition}

\begin{proof}
The family of entire fibers $\operatorname{St}^{\mathrm{add}}(F(r))$ is a stable seeded sub-zoo, so the defining class is nonempty. Arbitrary intersections of full stable subcategories are again full stable subcategories, and invariance under the designated exact functors is preserved under intersection. Hence the intersection is the least stable seeded sub-zoo.
\end{proof}

\begin{theorem}[Automatic residual promotion]\label{thm:auto-promotion}
Let $\mathcal Z$ be the minimal stable zoo. If $X\in\mathcal Z_s$, then every mixed residual
\[
\Delta^A_{s,r}(X)
\]
belongs to $\mathcal Z_r$.
\end{theorem}

\begin{proof}
By invariance,
\[
\widehat A_r\widehat C_{s\to r}X
\quad\text{and}\quad
\widehat C_{s\to r}\widehat A_sX
\]
are objects of $\mathcal Z_r$. Since $\mathcal Z_r$ is stable, it is closed under cofibers. Therefore the cofiber of the exchange morphism between these two objects, namely $\Delta^A_{s,r}(X)$, also lies in $\mathcal Z_r$.
\end{proof}

\status{proved}

\begin{corollary}[Redundancy of residual promotion in the stable calculus]\label{cor:promotion-redundant}
In the additive/stable two-axis calculus, ``promote every residual to a seed'' need not be imposed as an independent closure axiom. Once the zoo is stable and closed under transformation and coarsening, mixed residuals are already forced to lie in it.
\end{corollary}

This is a substantive simplification of the initial axiomatization. Residual promotion remains meaningful operationally for search and provenance, but mathematically it is redundant in the stable closure.

\subsection{What has and has not been proved}

Combining \cref{thm:free-double,thm:pointwise-stable,thm:residual-preservation,thm:auto-promotion} gives the following restricted universal construction.

\begin{theorem}[Additive/stable seeded two-axis completion]\label{thm:additive-stable-completion}
For a small additive indexed two-axis system with coherent additive transformation/exchange data and designated seeds, there is a canonical stable completion with the following properties:
\begin{enumerate}
    \item the resolution diagram extends universally from additive fibers to stable fibers;
    \item all coarsening and transformation functors become exact;
    \item every exchange morphism has a canonical mixed residual given by a cofiber;
    \item exact realizations preserve these residuals;
    \item the least stable seed zoo exists and is automatically closed under mixed residual formation.
\end{enumerate}
The stable completion has the adjunction universal property for stable resolution systems receiving the original additive resolution diagram; with compatible exchange data, residual preservation is automatic for exact realizations.
\end{theorem}

\begin{proof}
Items (1)--(2) are \cref{thm:pointwise-stable}; item (3) is \cref{def:universal-residual}; item (4) is \cref{thm:residual-preservation}; item (5) is \cref{prop:min-stable-zoo,thm:auto-promotion}. Initiality is the adjunction of \cref{thm:pointwise-stable}.
\end{proof}

\status{proved under the explicit additive/stable hypotheses}

The stronger nonadditive claim remains open. In particular, this theorem does not construct a free \emph{arbitrary} residual-complete double category in which residuals satisfy a universal cofiber property without assuming an additive/stable target. It does, however, close the universal-construction problem for the regime in which the present mixed residual is actually defined as a cofiber.


\subsection{Testing residual vanishing on orbit generators}

The universal completion makes residuals available, but an application still needs a way to prove they vanish. Exactness gives a reduction to generators, with an important qualification: operation-orbit generators, rather than the original seed list alone, generally control the invariant zoo.

\begin{lemma}[Residual formation is exact in its input]\label{lem:exact-residual-functor}
Let $P,Q:\mathcal C\to\mathcal D$ be exact functors between small stable $\infty$-categories, and let $\eta:P\Rightarrow Q$ be a natural transformation. Then
\[
\Delta_\eta(X):=\Cofib(\eta_X)
\]
is an exact functor $\mathcal C\to\mathcal D$. If $q:\mathcal D\to\mathcal O$ is exact, the subcategory
\[
\{X:q\Delta_\eta(X)\simeq0\}
\]
is thick.
\end{lemma}

\begin{proof}
Cofibers in the functor category are computed pointwise. The full subcategory of exact functors is closed under these finite colimits: a finite colimit of finite-colimit-preserving functors preserves finite colimits, and its value at zero is zero. Hence $\Delta_\eta$ is exact. The composite $q\Delta_\eta$ is exact as well. Its kernel is closed under shifts, cofibers, and retracts, which proves the last assertion.
\end{proof}

\begin{definition}[Operation orbit of the seeds]
In a stable indexed system, let $\mathscr P$ be the designated exact transformations and coarsenings. For each $r$, let $\mathcal G_r$ contain every image at level $r$ of a seed under a finite well-typed word in $\mathscr P$, including the empty word. Define
\[
\mathcal Z_r^{\mathrm{thick}}:=\operatorname{Thick}(\mathcal G_r).
\]
Refinement is included here only when an actual exact refinement functor has been separately specified. A multivalued lift relation is not silently treated as an exact functor.
\end{definition}

\begin{proposition}[Orbit generation of the invariant thick zoo]\label{prop:orbit-zoo}
The family $\{\mathcal Z_r^{\mathrm{thick}}\}_r$ is the least family of thick subcategories containing the seeds and preserved by all operations in $\mathscr P$.
\end{proposition}

\begin{proof}
Every operation sends an orbit generator to an orbit generator at the target level. Exactness and preservation of retracts therefore imply that it sends the thick closure of the source orbit into the thick closure of the target orbit. Thus the displayed family is invariant. Any invariant thick family containing the seeds contains all finite operation words applied to them, hence contains all $\mathcal G_r$ and their thick closures. This proves minimality.
\end{proof}

\begin{theorem}[Orbit-generator test for an exchange obstruction]\label{thm:generator-test}
Fix an exchange $\eta^A_{s,r}$ between exact functors, and an exact quotient $q_r$ of the target fiber. If
\[
q_r\Delta^A_{s,r}(g)\simeq0
\quad\text{for every }g\in\mathcal G_s,
\]
then
\[
q_r\Delta^A_{s,r}(X)\simeq0
\quad\text{for every }X\in\mathcal Z_s^{\mathrm{thick}}.
\]
The same assertion holds for $q_r=\id$, in which case the exchange is an equivalence on the entire generated thick zoo.
\end{theorem}

\begin{proof}
By \cref{lem:exact-residual-functor}, the objects on which $q_r\Delta^A_{s,r}$ vanishes form a thick subcategory. By hypothesis it contains $\mathcal G_s$, hence it contains $\operatorname{Thick}(\mathcal G_s)$. Vanishing of the unquotiented cofiber is equivalent to invertibility of the exchange map.
\end{proof}

\status{proved by exactness and thick generation}

This reduces a universal statement on a thick zoo to a statement on its operation orbit. It is not automatically a finite algorithm. The orbit may be infinite, and checking only the initial seeds does not suffice unless a separate argument propagates vanishing along every operation. In particular, this theorem does not turn finitely many numerical checks into proof of a universal arithmetic comparison.


\section{Research program and theorem ledger}

\begin{longtable}{>{\raggedright\arraybackslash}p{0.09\textwidth} >{\raggedright\arraybackslash}p{0.47\textwidth} >{\raggedright\arraybackslash}p{0.32\textwidth}}
\toprule
ID & Target & Current status\\
\midrule
\endhead
K1 & Resolution monotonicity and nested equivalence relations & Proved in \cref{thm:resolution-monotonicity}\\
K2 & Grothendieck total category & Proved in \cref{thm:total-category}\\
K3 & Two-axis factorization & Proved in \cref{thm:factorization}\\
K4 & Resolution ultrametric & Proved with a common state set, universal zeroth relation, nesting, and separation.\\
K5 & Least seed zoo / closure existence & Proved in \cref{thm:zoo-existence}\\
K6 & Finitary iterative zoo generation & Proved under continuity hypothesis in \cref{prop:finitary}\\
K7 & Lifting coherent transformation families to the total category & Proved in \cref{thm:lift-total}\\
K8 & Mixed residual vanishing criterion & Proved in stable fibers in \cref{prop:residual-vanish}\\
K9 & Exact two-axis normal form & Proved under local confluence in \cref{thm:normal-form}\\
K10 & Split residual expansion & Corrected source remainder shift and finite principal-path scope in \cref{thm:split-expansion}.\\
K11 & Stable residual filtration without splitness & Proved for a chosen path. General ordered independence is refuted by K24.\\
K12 & Stable obstruction quotient & Proved for the thick kernel. Constructor descent needs K28 or the directed hypothesis K29.\\
K13 & Free seeded two-axis double syntax & Term-construction proof with typed labels in \cref{thm:free-double}.\\
K14 & Universal residual-complete two-axis construction & Proved with additive/stable infinity-category and coherent exchange data. Nonadditive form open.\\
K15 & Effective reconstruction theorem for a nontrivial arithmetic instance & Lattice case proved in K51; the intended reciprocity instantiation remains open.\\
K16 & Descent transfer across a correspondence & Abstract lemma proved in \cref{thm:descent-transfer}\\
K17 & Seeded local-global reciprocity instantiation & Program theorem, arithmetic hypotheses open\\
K18 & Lean formalization of K1--K10 & Not yet implemented in this note\\
K19 & Pointwise stable completion universal property & Proved in \cref{thm:pointwise-stable}\\
K20 & Automatic residual promotion in stable zoos & Proved in \cref{thm:auto-promotion}\\
K21 & Functorial transport of obstruction quotients & Proved in \cref{thm:obstruction-transport}\\
K22 & Nonstable residual/repair obstruction test bed & Proved for abelian torsors in \cref{thm:torsor-obstruction}\\
K23 & Localization of total obstruction & Proved pathwise in \cref{thm:obstruction-localization}.\\
K24 & Ordered filtration independence & General claim refuted in \cref{prop:filtration-no-go,ex:typed-diamond}.\\
K25 & Semantic diamond criterion & Proved in \cref{thm:semantic-diamond}; total cofibers equivalent, not uniquely identified without more data.\\
K26 & Pure exchange event poset & States are ideals and histories are linear extensions, \cref{thm:event-ideals}.\\
K27 & Ideal-indexed residual diagram & Proved for coherent input; K32 constructs strict chain-natural instances.\\
K28 & Constructor-invariant obstruction quotient & Proved in \cref{thm:invariant-quotient}; can be zero, \cref{ex:quotient-collapse}.\\
K29 & Directed operational quotient & Killability is thick and constructor-invariant, \cref{thm:directed-invariant}.\\
K30 & Orbit-generator test for residual vanishing & Proved in \cref{thm:generator-test}; orbit need not be finite.\\
K31 & Noncollapse certificate & Exact equivariant detector criterion, \cref{prop:noncollapse}.\\
K32 & Strict semantic event realization & Proved for typed pure swaps and strict chain-natural cells, \cref{thm:strict-event-realization}.\\
K33 & Interaction residual and excision & Total-cofiber criterion, \cref{thm:interaction-cofiber}.\\
K34 & Context-independent event residuals & Equivalent to vanishing square interactions, \cref{thm:event-context-invariance}.\\
K35 & Principal-ideal reconstruction & Reduced excisive diagrams classified by diagrams on $P$, \cref{thm:principal-reconstruction}.\\
K36 & Universal excisive quotient & Constructor-compatible minimal thick kernel, \cref{thm:excisive-quotient}.\\
K37 & Tensor-cube residual formula & Proved by finite-colimit interchange and exact tensor, \cref{thm:tensor-cube}.\\
K38 & Integral gcd interaction formula & Explicit Koszul change of basis and cohomology, \cref{thm:integral-gcd}.\\
K39 & Prime-overlap kernel & Exact kernel and noncollapse for positive scalar labels, \cref{thm:prime-overlap-quotient}.\\
K40 & Finite-signature arithmetic certificate & All finite event diagrams in the stated scalar grammar, \cref{cor:finite-signature}.\\
K41 & Integral matrix collapse criterion & Verified Smith data classify free collapse versus torsion support, \cref{thm:matrix-collapse}.\\
K42 & Nonidentity arithmetic resolution tower & Derived reduction is exact and conservative, not faithful, \cref{thm:adic-tower}.\\
K43 & Coherent weighted exchange cells & Explicit source of residuals and base-change transport, \cref{thm:weighted-adic-exchange}.\\
K44 & Prime-power interaction cohomology & Full finite-level Koszul cohomology, \cref{thm:prime-power-koszul}.\\
K45 & Nilpotent excision collapse & Nonunit pairs force the prime-power quotient to zero, \cref{thm:nilpotent-collapse}.\\
K46 & Finite prime-component quotient & Exact quotient equivalence and tower coherence, \cref{thm:finite-prime-quotient}.\\
K47 & Based tower reconstruction & Uniform bounded ranks and compatible matrices required, \cref{thm:based-tower-reconstruction}.\\
K48 & Transient negative residual classes & Nonzero finite groups can form a pro-zero system, \cref{prop:transient-koszul}.\\
K49 & Completion--repair comparison defect & Explicit generic-fiber formula and vanishing criterion, \cref{thm:completion-repair-defect}.\\
K50 & Arithmetic fracture reconstruction & Perfect integral complexes and their residuals, \cref{thm:arithmetic-fracture,cor:residual-fracture}.\\
K51 & Effective arithmetic lattice gluing & Existence and uniqueness from specified generic identifications, \cref{thm:lattice-gluing}.\\
K52 & Compatible operator gluing & Preserves supplied integral polynomial relations, \cref{cor:lattice-operators}.\\
\bottomrule
\end{longtable}

\section{Immediate next formal steps}

The general ordered-filtration independence target remains refuted. A nonidentity arithmetic tower and an effective lattice-gluing instance now have proofs. Neither the finite prime-component quotient nor this lattice existence theorem supplies automorphic globalization.

\begin{enumerate}
\item \textbf{Beyond strict pure exchange.} Extend K32 to a specified non-strict enhanced realization or a rewrite signature with additional critical overlaps. Identify the missing coherent cells explicitly rather than assume that homotopy-category diamonds supply them.
\item \textbf{Beyond weighted scalar cells.} The derived coefficient tower is now constructed. Find a specified arithmetic functor whose natural comparison with coarsening, rather than an imposed scalar weight, has a nonzero residual. Determine its information-preserving repair or gluing mechanism.
\item \textbf{Arithmetic existence beyond lattices.} K51 closes a concrete lattice case of K15. For the intended K17 application, construct the generic arithmetic or automorphic data and the required identifications, rather than importing them as gluing hypotheses. A globally glued lattice with supplied operators is not an automorphic representation.
\item \textbf{Broader coefficient rings.} Determine which parts of the prime-support kernel and matrix criterion persist over a specified Dedekind ring or a higher-dimensional ring. A replacement for the integral Smith decomposition must be proved, not silently imported.
\item \textbf{Lean core.} Formalize the finite event-poset correspondence, strict naturality argument, closure statements, and integral certificate checks. This manuscript still has no Lean proof package.
\item \textbf{Beyond stability.} Develop another precisely typed nonadditive obstruction example beyond torsor sections. Do not identify a nonabelian obstruction set with a cofiber object without a comparison theorem.
\end{enumerate}

\section{Brainstorm: stronger directions}

The formalism above suggests several deeper lines of investigation.

\subsection{Resolution completion}
A resolution tower may admit an inverse limit
\[
X_\infty\simeq\varprojlim_r X_r.
\]
One can ask whether the seed zoo commutes with completion:
\[
\Zoo(X_\infty)
\stackrel{?}{\simeq}
\varprojlim_r\Zoo(X_r).
\]
This is only a proposed comparison. A common category and an actual comparison map must first be constructed before a residual is defined.

\subsection{Transformation completion}
Repeated transformations may generate an orbit category. Completing this orbit under limits, colimits, idempotent splitting, or derived closure may produce a canonical transformation hull.

\subsection{Bi-completion}
There are then two completion operations:
\[
\widehat{(-)}_T,
\qquad
\widehat{(-)}_R.
\]
After both completions and the comparison have been defined in a common category, the mixed question is whether
\[
\widehat{\widehat X_T}_R
\simeq
\widehat{\widehat X_R}_T.
\]
Given the specified comparison, its defect is a candidate completion-level mixed residual in a setting where such defects are defined. The comparison is now explicit in \cref{thm:completion-repair-defect}: even compatible finite-level repairs can fail to commute with completion, with a nonzero generic-fiber defect. An unconditional bi-completion claim is therefore false.

\subsection{Resolution spectrum}
For a fixed $X\in\Tcat_s$ and a specified detection predicate $P_r$ at each level $r\le s$, define:
\[
\operatorname{Spec}_R(P;X)
=
\{r\le s:P_r(C_{s\to r}X)\}.
\]
This could behave like a support or spectrum for distinguishability.

\subsection{Residual genealogy}
Residual promotion creates a directed graph
\[
s\to\Delta_1\to\Delta_2\to\cdots.
\]
Classify whether genealogies terminate, become periodic, stabilize, or grow indefinitely. A terminating genealogy would correspond to finite obstruction resolution. Periodic genealogy may encode a persistent but structured defect.

\subsection{Seed minimality}
Given a generated zoo $\Zoo$, ask for the smallest seed set $\Seeds$ with
\[
\Zoo(\Seeds)=\Zoo.
\]
This is a mathematical compression problem. Nonuniqueness of minimal seeds could itself be meaningful.

\subsection{Dual zoos}
If the calculus carries a duality $D$, study
\[
D\Zoo(\Seeds)
\quad\text{versus}\quad
\Zoo(D\Seeds).
\]
The residual of dualization against zoo generation measures whether the generative grammar is duality-complete.

\subsection{Proof-resolution}
A theorem itself may have resolutions. At coarse resolution one stores only a proposition and certificate. At finer resolution one exposes lemmas, dependencies, computational witnesses, and failed branches. This creates a resolution geometry on proofs and may be useful for machine reasoning and proof compression.

\subsection{World-resolution}
A foundational extension can assign resolution to models or ``worlds'' themselves. Morphisms between worlds then have two axes: change of world and change of resolution within a world. This may connect the current calculus to the earlier compatibility/residual/obstruction maps between models of a foundational theory.

\section{Provisional philosophical interpretation}

The mathematics does not require the following interpretation, but the two-axis formalism sharpens it.

A structure may be characterized simultaneously by:
\begin{itemize}
    \item what transformations it admits;
    \item what resolutions distinguish its internal states;
    \item what survives both kinds of movement;
    \item what residuals appear when the two directions fail to commute;
    \item what obstructions prevent refinement, transformation, or global reconstruction.
\end{itemize}

This suggests the strengthened ontological phrase:
\[
\boxed{
\text{to remain identifiable is to remain coherent under both transformation and resolution.}
}
\]
The mathematical content of the project is not this sentence itself. The mathematical content is the attempt to make ``coherent,'' ``transformation,'' ``resolution,'' ``residual,'' and ``reconstruction'' precise enough that such a sentence can be tested rather than merely asserted.

\section{Conclusion}

The first formalization supports a concrete picture:
\[
\boxed{
\begin{array}{c}
\text{seeds generate states},\\
\text{transformations generate representations},\\
\text{resolutions generate distinguishability},\\
\text{exchange maps relate the axes},\\
\text{residuals measure failed exchange},\\
\text{residuals become new seeds},\\
\text{obstructions classify persistent failure},\\
\text{descent reconstructs compatible local structure}.
\end{array}
}
\]
The retained categorical core supplies explicit constructions under fixed hypotheses. Stable completion provides functorial cofibers, exact realizations preserve residuals, and the invariant stable zoo contains them automatically. Each normalization path carries a residual filtration, but the filtration is genuinely history-sensitive. The corrected intrinsic object is an ideal-indexed residual diagram for a specified coherent exchange system, whose terminal residual is independent of schedule.

Likewise, a thick obstruction quotient is not automatically a new environment for the original constructors. Their action on the quotient requires an invariant kernel. The invariant saturation gives the universal repair, but its possible collapse is a substantive constraint on admissibility. Directed common postcomposition makes operational killability itself invariant and thick, while exact equivariant detectors certify noncollapse. Finally, thick orbit generation reduces exchange vanishing to a well-defined generator test without promising that the test is finite.

The next layer is now explicit. Strict chain-natural cells construct semantic exchange diagrams. Their second cofibers measure excision, and vanishing of those interaction objects makes event residuals independent of unrelated prior events. Principal-ideal reconstruction retains exactly the poset diagram needed to assemble the whole residual history. In the integral models, common prime support identifies what must be removed, while free cofiber components identify a genuine collapse risk. These are concrete semantic and arithmetic results, not a conclusion drawn solely from the breadth of the vocabulary.

The arithmetic tower sharpens that boundary further. Finite resolution can turn a nonzero nilpotent interaction into a generator of the whole perfect category. Several prime components permit nontrivial excisive quotients, but deleting a component before completion loses a generic fiber. Retaining the localized and completed branches with their overlap instead gives a proved arithmetic fracture square and an effective lattice-gluing theorem. These use established algebraic mechanisms to make the calculus concrete.

The results do not establish a new foundation or a reciprocity theorem. The remaining arithmetic content requires actual generic parameters, compatible actions, and the specific existence theorem demanded by the intended correspondence. Neither a torsor trivialization, a reconstructed event diagram, nor a glued integral lattice is an automorphic realization by definition.


\begin{thebibliography}{99}
\addcontentsline{toc}{section}{References}
\bibitem{FairbanksShulman2026}
A.~D. Fairbanks and M.~Shulman,
\emph{Doubly Weak Double Categories},
arXiv:2506.23651. Background reference for double computads.
The paper is background for double computads. The present free strict syntax result is proved directly by terms and relations in \cref{thm:free-double}; no unchecked theorem-number identification is needed.

\bibitem{SaunierWinges2026}
V.~Saunier and C.~Winges,
\emph{On exact categories and their stable envelopes},
arXiv:2502.03408v2, 28 November 2025.
Definition 3.1 and the following universal property give the free stable envelope $\operatorname{St}^{\mathrm{add}}$ and its adjunction with the forgetful functor from stable to additive categories.
\url{https://arxiv.org/abs/2502.03408v2}.

\bibitem{StacksVerdier}
The Stacks Project Authors,
\emph{Quotients of triangulated categories}, Section 13.6, Tag 05RA, especially Lemmas 13.6.2, 13.6.8, and 13.6.9, accessed 26 September 2026.
The quotient construction, kernel description, and exact universal property are used in \cref{thm:obstruction-quotient,thm:obstruction-transport}; Section 13.3, Tag 05QK, records the octahedral axiom TR4 underlying the homotopy-category form of \cref{thm:stable-filtration}. Functorial cofibers require the stable enhancement fixed in the text.
\bibitem{Drew2015}
B.~Drew,
\emph{Verdier quotients of stable quasi-categories are localizations},
arXiv:1511.08287v1, 26 November 2015, Definition 1.1 and Theorem 1.3.
\url{https://arxiv.org/abs/1511.08287v1}.

\bibitem{StacksTorsors}
The Stacks Project Authors,
\emph{First cohomology and torsors}, Section 21.4, Tag 03AG, accessed 26 September 2026.

\bibitem{StacksCones}
The Stacks Project Authors,
\emph{Cones and termwise split sequences}, Section 13.9, Tag 014D,
Definition 13.9.1 and Lemma 13.9.2, checked 26 September 2026.
\url{https://stacks.math.columbia.edu/tag/014D}.

\bibitem{StacksKoszul}
The Stacks Project Authors,
\emph{The Koszul complex}, Section 15.29, Tag 0621,
Definitions 15.29.1--2 and Lemmas 15.29.4, 15.29.7--8, 15.29.12,
checked 26 September 2026.
The cited section uses homological grading. This manuscript puts the exterior
$k$-th term in cohomological degree $-k$ and states all differentials explicitly.
\url{https://stacks.math.columbia.edu/tag/0621}.

\bibitem{StacksPerfect}
The Stacks Project Authors,
\emph{Perfect complexes}, Section 15.76, Tag 0656,
Definition 15.76.1 and Lemma 15.76.9, checked 27 September 2026.
\url{https://stacks.math.columbia.edu/tag/0656}.

\bibitem{StacksDerivedTensor}
The Stacks Project Authors,
\emph{Derived tensor product}, Section 15.60, Tag 06XY,
checked 27 September 2026. K-flat models and the exact derived tensor construction.
\url{https://stacks.math.columbia.edu/tag/06XY}.

\bibitem{StacksRlim}
The Stacks Project Authors,
\emph{Rlim of abelian groups}, Section 15.88, Tag 07KV,
checked 27 September 2026. The product model and the Mittag--Leffler case of derived inverse limits.
\url{https://stacks.math.columbia.edu/tag/07KV}.

\bibitem{StacksLocalization}
The Stacks Project Authors,
\emph{Localization}, Section 10.9, Tag 00CM,
checked 27 September 2026. The universal property and exactness on modules.
\url{https://stacks.math.columbia.edu/tag/00CM}.

\bibitem{StacksDerivedCompletion}
The Stacks Project Authors,
\emph{Derived Completion}, Section 15.93, Tag 091N,
Lemma 15.93.12, checked 27 September 2026. Inverting a defining element before derived completion gives zero.
\url{https://stacks.math.columbia.edu/tag/091N}.

\end{thebibliography}

\appendix
\section{Compact definition sheet}

\begin{align*}
\text{Resolution system:}\quad
&F:\Rcat^{op}\to\Cat.\\[0.4em]
\text{Fiber:}\quad
&\Tcat_r=F(r).\\[0.4em]
\text{Coarsening:}\quad
&C_{s\to r}:\Tcat_s\to\Tcat_r.\\[0.4em]
\text{Total category:}\quad
&\K=\int_{\Rcat}F.\\[0.4em]
\text{Resolution equivalence:}\quad
&X\sim_rY\iff C_{t\to r}X\cong C_{t\to r}Y.\\[0.4em]
\text{Lift set:}\quad
&\Lift_{r\to s}(X)=\{Y:C_{s\to r}Y\cong X\}/\cong.\\[0.4em]
\text{Exchange map:}\quad
&\eta^A_{s,r}:A_rC_{s\to r}\Rightarrow C_{s\to r}A_s.\\[0.4em]
\text{Stable residual:}\quad
&\Delta^A_{s,r}(X)=\Cofib(\eta^A_{s,r,X}).\\[0.4em]
\text{Stable obstruction quotient:}\quad
&\mathcal O_{\mathscr M}=\mathcal Z/\operatorname{Thick}(\mathcal K_{\mathscr M}).\\[0.4em]
\text{Invariant obstruction quotient:}\quad
&\mathcal O^{\mathrm{inv}}_{\mathscr M}=\mathcal Z/\mathcal J_{\mathscr M}.\\[0.4em]
\text{Event residual diagram:}\quad
&\mathcal F_X(I)=\Cofib(X_{\varnothing}\to X_I).\\[0.4em]
\text{Obstruction class:}\quad
&\Obs_{\mathscr M}(X)=q(X).\\[0.4em]
\text{Seed zoo:}\quad
&\Zoo(\Seeds_0)=\mu A.\bigl(\Seeds_0\cup\Phi(A)\bigr).\\[0.4em]
\text{Resolution ultrametric:}\quad
&d_R(X,Y)=2^{-\nu(X,Y)},\qquad 2^{-\infty}=0.\\[0.4em]
\text{Reconstruction target:}\quad
&\Omega_G=0\iff\text{the specified descent datum is effective}.
\end{align*}


\section{Preserved audit and verification record for v0.6.0}\label{app:audit}

This section preserves the preceding release's audit. The present release's work is recorded in \cref{app:audit-v070}.

\subsection{Coverage and evidence}
The v0.5.0 LaTeX source was read in full. This release preserves the running manuscript and revises its statements and proofs rather than replacing it with a shortened summary. The review is an authoring-session self-review, not an independent referee report. The exact source snapshot, change inventory, executable tests, and output are included in the companion package.

The primary audit verdict for the former package is \emph{correct only after stated restrictions and corrections}. In particular, the ordered-filtration ambition fails in general, whereas the retained path-relative theorem survives. Stable quotient existence survives, but it alone does not prove constructor descent or noncollapse. No Lean verification is claimed.

\subsection{Mathematical corrections}
\begin{enumerate}
\item \textbf{Resolution:} fixed a common state set and a universal zeroth equivalence relation in the ultrametric construction. Distinguished isomorphism classes of lifts from framed lift groupoids. Excluded $\log 0$ from ordinary branching information.
\item \textbf{Exchange:} distinguished invertible comparison maps from exact functors. Required complete local exchange rules for the syntactic normal-form statement. Replaced the split cofiber identification by the correct shift $\Delta\simeq E[1]$. Removed an unsupported recursive termination claim for residual branches.
\item \textbf{Coherence:} separated a normal object, its comparison map, and a filtration. Noninvertible exchange cells are extended using the full functor-category stable-envelope universal property, not just a mapping-space adjunction.
\item \textbf{Quotients:} distinguished killability, thick killability, and constructor-invariant thick killability. Stated universal uniqueness in the enhanced categorical sense. Corrected the postcomposition orientation in the directed condition.
\item \textbf{Descent and scope:} included all overlap degrees and descent of morphisms. Required reflection as well as preservation of obstruction vanishing. Corrected the quotient functor's type. Kept arithmetic globalization and unrestricted reciprocity outside the proved scope.
\item \textbf{Derived quantities:} replaced an identity-dominated commutation-depth formula and typed the branching profile along an actual lift chain. Marked zoo growth and search costs as presentation-dependent. No minimum or decision procedure is asserted by an optimization display.
\end{enumerate}

\subsection{Exact finite regression}
The companion script \texttt{checks/check\_models.py} uses integer matrices, rational Gaussian elimination, and finite labelled words. It contains no floating-point tests. The recorded run uses Python 3.13.5 and checks every two-sort word of length $0$ through $10$, including the empty word and words containing only one sort.

\begin{center}
\begin{tabular}{lr}
\toprule
Quantity & Checked count\\
\midrule
Initial words & 2,047\\
Reachable word states, summed over initial words & 82,499\\
Elementary rewrite edges, summed over initial words & 168,454\\
Largest event poset & 25 events\\
Products in the three-constructor example & 9\\
Binary supports in the three-coordinate example & 8\\
Exact rank checks for cone models & 6\\
\bottomrule
\end{tabular}
\end{center}

The word test constructs reachable labelled words and order ideals independently, then compares their states and edges. The constructor test confirms the non-thick killability and invariant-collapse examples, together with a nonzero directed projection model. The cone tests compare the ordered dimensions $(0,1,3)$ and $(0,2,3)$, verify the split projection's shifted cofiber, and distinguish the identity and zero maps with equal endpoints.

All these finite checks passed. They test the stated finite models and implementation paths. The universal statements rest on their written proofs, not on extrapolation from the finite enumeration. The package's result record stores the exact script hash and bounds. Build and visual-review results are recorded separately from the mathematical test status.

\subsection{Source verification boundary}
The stable-envelope input was checked in Saunier--Winges, arXiv:2502.03408v2, Definition 3.1 and the subsequent functor-category universal property. The stable localization input was checked in Drew, arXiv:1511.08287v1, Definition 1.1 and Theorem 1.3. The quotient kernel and triangulated formulations were checked in the cited Stacks sections. The free strict syntax theorem now has its own term-construction proof. This was an exact-dependency check, not a comprehensive novelty search.

The manuscript does not certify a new axiom of set theory, a classification of all mathematical structures, a physical law, or a Langlands reciprocity theorem. Its new claims concern the explicitly defined exchange grammar and stable constructor systems.

\clearpage
\section{Audit and verification record for v0.7.0}\label{app:audit-v070}

\subsection{Scope, proof status, and provenance}
The source base is the supplied v0.6.0 manuscript and its complete package. This release preserves that manuscript, integrates new proof blocks, and updates the abstract, theorem ledger, and next obligations. It does not claim a fresh independent audit of every retained earlier theorem. The new proofs received an authoring-session self-review. No independent referee review or Lean compilation was performed.

The three executed routes were strict semantic realization, interaction/excision reconstruction, and integral quotient classification. They produced K32--K41 under their written hypotheses. The arithmetic test bed uses perfect complexes over $\mathbb Z$ and explicit finite integer data. It does not discharge a Langlands-facing existence hypothesis.

\subsection{Adversarial boundaries}
\begin{enumerate}
\item Strict naturality supplies the event diagram only for the specified pure grammar. Weak cells and additional rewrite relations retain separate coherence obligations.
\item A square's commutativity is not its excision. The second cofiber vanishes precisely for a homotopy pushout, not for every commuting square.
\item Context-independent event residuals do not identify ordered filtrations. Nor do they split extensions along dependent events: the $\mathbb Z/p^2$ chain is an explicit counterexample.
\item The total normalization residual $C(N)$ and a cube's total cofiber $K(n_1,\ldots,n_m)$ are different objects. The $(6,10,15)$ example refutes an inference from highest interaction vanishing to all lower interactions vanishing.
\item The integral noncollapse theorem uses positive scalar labels. Two zero labels give an interaction with a free unit summand and collapse the entire quotient. The matrix criterion identifies the analogous free-cofiber failure.
\item Equality of the kernel with that of localization is proved. No unproved equivalence of complete quotient categories or idempotent completions is substituted for it.
\end{enumerate}

\subsection{Exact finite checks}
Both regression scripts ran successfully on Python 3.13.5. They use exact integer identities and finite-field elimination, with no random sampling or numerical tolerance. The inherited check was rerun with every two-sort word through length 10: 2,047 initial words, 82,499 states, and 168,454 edges. Its constructor and filtration counterexamples also passed.

The new check covers the following explicitly bounded populations.
\begin{center}
\begin{tabular}{p{0.72\textwidth}r}
\toprule
Check population & Count\\
\midrule
Positive label vectors, arity 1--4, entries 1--6 & 1,554\\
Integral exterior-power chain-map identities & 5,910\\
Integral differential-square identities & 4,356\\
Cohomology dimensions over $\mathbb F_2,\mathbb F_3,\mathbb F_5,\mathbb F_7$ & 29,856\\
Initial scalar two-sort words through length 8 & 511\\
Ideal states in those scalar diagrams & 6,917\\
Cover arrows and eligible event contexts & 10,660\\
Elementary independent-event squares & 4,931\\
Explicit localized context-map inverses & 5,019\\
Ideal-pair combinatorial excision checks, word length at most 6 & 5,289\\
Ordered pairs from an explicit 11-matrix test family & 121\\
Matrix-pair cohomology dimensions over five finite fields & 1,815\\
\bottomrule
\end{tabular}
\end{center}

The matrix family includes identity, zero, rank-deficient, triangular, and torsion examples. Its 121 pairs are exhaustive for that listed family, not for all integer matrices. Nine pairs have two free cone components. Additional signed/zero-vector checks, the excluded all-zero pair, and the named $(6,10)$, $(6,10,15)$, and nonsplit prime-chain examples are recorded separately. No primary case within the stated bounds was skipped.

The integral normal-form checks verify unimodular exterior-power chain isomorphisms. Finite-field ranks provide a distinct check on signs and degree dimensions. These are reproducible implementation checks. They neither instantiate every stable category nor replace the written universal proofs.

\subsection{Sources and remaining review}
The exact source check used Stacks Tags 014D, 0621, and 05RA, and Drew's arXiv:1511.08287v1, Definition 1.1 and Theorem 1.3. The source record separates the homological Koszul convention from the cohomological convention used here. Direct finite-colimit, lattice, and integer arguments are supplied for the new deductions. This was not a comprehensive novelty search. The companion audit contains the dependency map, script hashes, source record, build logs, and visual-review record.

Independent mathematical review remains outstanding. K15, K17, K18, arbitrary higher-coherent exchange signatures, and the nonadditive universal construction remain open within this project.

\clearpage

\clearpage
\section{Audit and verification record for v0.8.0}\label{app:audit-v080}

\subsection{Scope and proof boundary}
The supplied v0.7.0 PDF, source, and full package are preserved unchanged. Their hashes are recorded in the release. This continuation adds K42--K52, updates the restricted lattice status of K15, and retains the earlier open reciprocity obligations. The new proofs have received an authoring-session self-review, not an independent referee review or Lean certification. No fresh complete audit of all inherited theorems is claimed.

The three executed mathematical routes were derived coefficient towers, excision versus completion, and arithmetic reconstruction with retained overlaps. The first produces explicit finite-level interactions. The second refutes unconditional commutation of repair and completion and proves a generic-fiber criterion. The third proves a fracture square and effective lattice gluing, with compatible operator descent.

\subsection{Adversarial checks}
\begin{enumerate}
\item Derived reduction is exact, but naive tensor on arbitrary modules is not. The infinite periodic resolution over $\mathbb Z/p^2$ prevents replacing a perfect two-term cone by its degree-zero cohomology.
\item Canonical tensor/base-change exchanges have zero residual. Nonzero scalar cells are specified additional data, not an unexplained consequence of coarsening.
\item Prime-power interactions can generate the whole perfect category even for nonzero labels. A multi-prime quotient removes whole components and can still collapse if every prime is selected.
\item Coherent finite reductions do not imply commutation with an infinite inverse limit. The field factors in K49 exhibit the precise missing information.
\item The based tower theorem requires common bounded ranks and compatible matrices. Arbitrary incompatible level choices are not claimed to lift.
\item The fracture square reconstructs a supplied complex. K51 separately proves effective existence for independent lattice data, but requires the generic identifications explicitly. Operator gluing requires exact intertwining, not only matching characteristic polynomials.
\item Generic-fiber defects are computed in $D(\widehat B)$, not asserted to be perfect objects over the completed ring. Their cofiber is the $[1]$ shift of the stated fiber.
\end{enumerate}

\subsection{Reproducible finite checks}
All three check scripts passed in this session with Python 3.13.5. The inherited two scripts reproduce their earlier ranges. The new script additionally uses SymPy 1.14.0 for integer Hermite normal forms; each resulting lattice basis is checked by its exact finite residues and index. There is no random sampling or numerical tolerance.

The new suite checks all 1,080 ordered coefficient pairs in the eight listed rings $\mathbb Z/q$ with $q\in\{2,3,4,5,8,9,16,25\}$. It enumerates kernels and images to check three cohomology degrees and verifies middle-cohomology elementary divisors by torsion counts. It also checks 160 annihilator-transition tuples, 1,728 exchange coherences, 72 prime-component projections, 2,304 positive-label pair instances, all 30 pairs of subspaces of $\mathbb F_2^2$ and $\mathbb F_3^2$, 2,430 operator-preservation instances, 6,100 rank-one rational lattice membership tests, and 972 principal-part reconstruction witnesses. Exact bounds and nonclaims are in the JSON report.

These are finite regression checks, not computational proofs of universal categorical statements or assertions about every $p$-adic input. The limiting and gluing theorems have separate written proofs. The arithmetic-square argument and nilpotent-cone argument use classical mechanisms; no historical novelty is inferred from their packaging here.

\subsection{Build and source verification}
The standalone source is regenerated from the preserved base and the two new proof fragments by \texttt{audit/integrate\_v080.py}. The source diff, run reports, source checks, build log, and rendered-page audit accompany the release. The first general web search returned irrelevant results and was not used as evidence; the cited primary Stacks Project pages were then opened directly. No external source text or font files are bundled. The lack of a Lean executable in this runtime is recorded as an environment limitation, not a mathematical failure.

\section{Version ledger}

\begingroup\small\raggedright
\begin{description}[leftmargin=8em,style=nextline]
\item[v0.8.0] Adds nonidentity derived arithmetic coarsenings, coherent weighted exchanges, the prime-power Koszul formula, and nilpotent excision collapse. Classifies finite prime-component quotients and their coherence. Proves uniform based-tower reconstruction, computes transient negative classes, and identifies the generic-fiber defect of completion versus repair. Proves arithmetic fracture reconstruction, effective lattice gluing, and compatible operator descent. Preserves the v0.7.0 package and adds exact finite tests. Written proofs and self-review only.
\item[v0.7.0] Constructs strict chain-natural event diagrams. Defines interaction residuals and proves the total-cofiber and context-independence criteria. Proves principal-ideal reconstruction for reduced excisive event diagrams. Constructs the universal constructor-compatible excisive quotient. Gives tensor-cube, integral gcd, prime-overlap kernel, finite-signature, and integral matrix collapse theorems. Adds explicit nonsplitting, higher/lower interaction, and zero-label collapse boundaries, together with exact finite regression and a preserved full v0.6.0 package. Written proofs and self-review only.
\item[v0.6.0] Audits size, variance, categorical regimes, lift multiplicity, cofiber shifts, exactness, descent, and quotient interpretation. Refutes general ordered filtration independence, including inside the typed two-axis grammar. Proves the semantic diamond criterion and the order-ideal description of pure exchange histories. Constructs the coherent ideal-indexed residual diagram. Proves constructor-invariant quotient and directed invariance theorems, records a quotient-collapse counterexample, and adds exact orbit-generator tests and a noncollapse certificate. Includes finite exact regression checks. Written proofs only, not Lean verification.
\item[v0.5.0] Combines the stable residual filtration with the Verdier obstruction quotient to prove obstruction localization along any chosen normalization path: if all elementary exchange residuals vanish in the obstruction quotient, so does the total normalization residual; contrapositively, a persistent total obstruction forces at least one persistent elementary exchange obstruction. Under directed annihilators the statement upgrades from quotient-vanishing to operational killability.
\item[v0.4.0] Adds a direct nonstable obstruction test bed: local sections of an abelian torsor produce overlap residuals, local changes of trivialization act as repairs, and the resulting \v{C}ech $H^1$ class vanishes exactly when the local sections can be repaired to reconstruct a global section. This closes the first descent test-bed obligation and shows that the residual/repair/obstruction pattern recovers a classical cohomological mechanism.
\item[v0.3.0] Proves the stable residual filtration along a chosen normalization path by iterated octahedral cofiber sequences. Constructs the universal stable obstruction quotient by the thick closure of operationally killable residuals, proves exact detection under a directed common-annihilator hypothesis, and proves functorial transport of obstruction quotients under exact correspondences. K11 and K12 are closed in these stable forms; path-independent filtration and nonstable obstruction remain open.
\item[v0.2.0] Adds the free seeded two-axis double-computad universal property and proves the pointwise free stable completion theorem for additive indexed systems. Mixed residuals become canonical cofibers, exact realizations preserve them, and residual promotion is shown to be automatic in the minimal stable zoo. K13 is closed at the strict syntactic level and K14 is closed under additive/stable hypotheses; the nonadditive residual universal construction remains open.
\item[v0.1.0] Initial continuous formalization. Introduces indexed resolution categories, total Grothendieck category, factorization, resolution ultrametric, seed zoo closure, exchange transformations, stable residuals, exact normal forms, split residual expansion, obstruction taxonomy, descent transfer, and a seeded local-global reciprocity blueprint.
\end{description}
\endgroup

\end{document}
